\documentclass[11pt]{article}
\usepackage[T1]{fontenc}
\usepackage[utf8]{inputenc}
\usepackage{lmodern}
\usepackage{amsmath,amssymb,amsthm,mathtools}
\usepackage[margin=1.05in]{geometry}
\usepackage{microtype,booktabs,array,enumitem}
\usepackage[dvipsnames]{xcolor}
\usepackage[colorlinks=true,linkcolor=blue!45!black,citecolor=blue!45!black,urlcolor=blue!45!black]{hyperref}
\usepackage[capitalise,nameinlink]{cleveref}
\hypersetup{pdftitle={The linear Paulsen bound: a streamlined proof},pdfauthor={}}

\newtheorem{theorem}{Theorem}[section]
\newtheorem{lemma}[theorem]{Lemma}
\newtheorem{proposition}[theorem]{Proposition}
\newtheorem{corollary}[theorem]{Corollary}
\theoremstyle{definition}
\newtheorem{definition}[theorem]{Definition}
\theoremstyle{remark}
\newtheorem{remark}[theorem]{Remark}

\newcommand{\R}{\mathbb R}
\newcommand{\E}{\mathbb E}
\newcommand{\Prob}{\mathbb P}
\newcommand{\one}{\mathbf 1}
\newcommand{\F}{\mathrm F}
\newcommand{\op}{\mathrm{op}}
\newcommand{\eps}{\varepsilon}
\newcommand{\ip}[2]{\langle #1,#2\rangle}
\newcommand{\norm}[1]{\lVert #1\rVert}
\newcommand{\fro}[1]{\lVert #1\rVert_{\F}}
\newcommand{\opn}[1]{\lVert #1\rVert_{\op}}
\newcommand{\Lap}{\mathcal L}
\DeclareMathOperator{\diag}{diag}
\DeclareMathOperator{\Diag}{Diag}
\DeclareMathOperator{\tr}{tr}
\DeclareMathOperator{\Cov}{Cov}
\DeclareMathOperator{\Var}{Var}
\DeclareMathOperator{\rank}{rank}
\DeclareMathOperator{\im}{im}
\numberwithin{equation}{section}
\setlength{\emergencystretch}{2em}
\allowdisplaybreaks[1]

\title{The linear Paulsen bound:\\ a streamlined proof}
\author{}
\date{October 2026}

\begin{document}
\maketitle
\vspace{-1.5em}
\begin{abstract}
Every real $\eps$-nearly equal-norm Parseval frame of $n$ vectors in $\R^d$
is within squared Frobenius distance $C\eps d$ of an equal-norm Parseval frame,
where $C$ is a universal constant. This is a reorganised and partly simplified
account of a recent proof of this bound (produced with an AI system and
formalised in Lean). The argument has one deterministic engine---a static
row-scaling theorem driven by an $\ell_\infty$ Poisson estimate for the
squared-Gram Laplacian---and two random ``seeds'' that make this engine
applicable: a filtered Gaussian tangent when $d\gtrsim\log n$, and a
row-tangent Gaussian when $n\gtrsim d^2$. Compared with the original
manuscript, a sharper remainder estimate lets the many-row seed work at every
error level below an absolute constant, which removes the random block
partition and the $d^6$ block size from the global argument; the scaling
lemmas are collected into one toolbox; and all probabilistic estimates are
proved in full.
\end{abstract}

\tableofcontents

\section{The result and the plan}\label{sec:intro}

\subsection{Statement}
A frame of $n$ vectors in $\R^d$ is a matrix $X\in\R^{n\times d}$ whose $i$th
row is $x_i^T$. Distances are \emph{squared} Frobenius distances
$\fro{X-Y}^2=\sum_i\norm{x_i-y_i}^2$, with no normalisation. Throughout,
\[
 a=d/n .
\]
$X$ is \emph{Parseval} if $X^TX=I_d$ and \emph{equal-norm} if
$\norm{x_i}^2=a$ for all $i$; an equal-norm Parseval frame is an \emph{ENP frame}.
$X$ is an \emph{$\eps$-nearly ENP frame} if
\begin{equation}\label{eq:nearly}
 (1-\eps)I_d\preceq X^TX\preceq(1+\eps)I_d,
 \qquad (1-\eps)a\le\norm{x_i}^2\le(1+\eps)a\quad(1\le i\le n).
\end{equation}

\begin{theorem}[Linear Paulsen bound]\label{thm:main}
There is an absolute constant $C$ such that for all $1\le d\le n$, all
$0<\eps<1$ and every real $\eps$-nearly ENP frame $X\in\R^{n\times d}$ there is
an ENP frame $W\in\R^{n\times d}$ with $\fro{X-W}^2\le C\eps d$.
\end{theorem}

The bound is optimal up to the constant. Kwok, Lau, Lee and
Ramachandran~\cite{KLLR} proved the first bound independent of $n$,
$O(\eps d^{13/2})$; Hamilton and Moitra~\cite{HM} improved it to
$O(\eps d^2)$; Lau and Ramachandran announced the linear bound~\cite{LR}. The
present text reorganises the proof of~\cite{Orig}; \cref{app:notes} lists what
was changed.

For a Parseval $U$ write $P=UU^T$ for the orthogonal projection onto its column
space; then $P_{ii}=\norm{u_i}^2$. The problem has an equivalent form in which
the error is additive.

\begin{theorem}[Projection form]\label{thm:projection}
There is an absolute constant $C$ such that for every $n\ge1$ and every
rank-$d$ orthogonal projection $P$ on $\R^n$, with
$\beta=\max_i|P_{ii}-d/n|$, there is a rank-$d$ orthogonal projection $Q$ with
$Q_{ii}=d/n$ for all $i$ and $\fro{P-Q}^2\le Cn\beta$.
\end{theorem}

Both $\beta$ and $\fro{P-Q}$ are unchanged when $(P,Q)$ is replaced by
$(I-P,I-Q)$, so in \cref{thm:projection} one may assume $d\le n/2$. The two
theorems are derived from each other and from the main proposition in
\cref{sec:assembly}.

\paragraph{Conventions.} $c,C$ denote positive absolute constants whose value
may change between occurrences; constants with subscripts or names are fixed
once chosen. $\opn{\cdot}$ is the Euclidean operator norm, $\one$ the all-ones
vector, $\Pi=I-n^{-1}\one\one^T$ the projection onto $\one^\perp$, $\Diag(x)$ the
diagonal matrix with diagonal $x$, and
$\diag(M)$ the diagonal of $M$ as a vector. $A\preceq B$ is the Loewner order.
A Laplacian inequality ``on $\one^\perp$'' means the corresponding quadratic
forms are compared on vectors orthogonal to $\one$. All Gaussian vectors are
real and centred.

\subsection{The plan}\label{subsec:plan}

\paragraph{The engine: scaling.}
Let $U$ be Parseval with projection $P$ and diagonal $p=\diag P$. For
$s\in\R^n$ the matrix $\Diag(e^s)U$ has a column space with projection $P(s)$,
and $s\mapsto\diag P(s)$ is the gradient of the convex function
$\frac12\log\det(U^T\Diag(e^{2s})U)$, whose Hessian at $s=0$ is $2L_P$, where
\[
 L_P=\Diag(p)-P\circ P
\]
is the graph Laplacian with edge weights $P_{ij}^2$ (\cref{lem:energy}). To
first order, then, changing the diagonal by $\delta\perp\one$ requires
$s\approx\frac12L_P^{+}\delta$, and the cost of the move is about
$s^TL_Ps\approx\frac14\delta^TL_P^{+}\delta$. If $L_P$ has spectral gap
$\lambda$ and $|\delta_i|\le\eps a$, the cost is at most $n\eps^2a^2/\lambda$.
For a ``random-looking'' $P$ one has $\lambda\approx a$ and the cost is
$\eps^2 d$, but $L_P$ can be arbitrarily badly connected: for an exact direct
sum the components cannot exchange diagonal mass at all.

\paragraph{Perturb, then scale.}
Perturb $U$ by a random horizontal tangent of size $t=\sqrt{K_0\eps}$, at cost
$O(t^2d)=O(\eps d)$. Generic entries of the new projection then have size
$t\sqrt{a/n}$, so most edges acquire weight $\gtrsim at^2/n$, and the random
walk with rates $P_{ij}^2$ reaches any half of the vertices in time $O(1/(at^2))$:
the \emph{barrier constant} (\cref{def:barrier,lem:core}) is $O(1/(at^2))$. The
static balancing theorem (\cref{thm:balancing}) then removes any diagonal error
$\beta$ with $\beta\cdot O(1/(at^2))\le\frac12$ exactly, at cost
$\frac{15}4n\beta=O(t^2d)$. The $\ell_\infty$ (rather than spectral) formulation is
essential: the scaling must stay in a box $\norm{s}_\infty\le1$, and an $\ell_2$
argument loses a factor $\sqrt n$.

\paragraph{The difficulty.}
So the perturbed frame must have diagonal error $\le\theta at^2$ in
$\ell_\infty$ for a small fixed $\theta$. The initial error $\eps a=at^2/K_0$ is
fine. But a random tangent of size $t$ moves each diagonal entry
\emph{linearly} by about $t\sqrt{a/n}$, which exceeds $at^2$ as soon as
$\eps\lesssim1/d$. Two different constructions avoid this.

\paragraph{Seed 1: many rows ($n\ge Bd^2$, \cref{sec:manyrow}).}
Perturb each row \emph{orthogonally to itself} and renormalise. Row norms stay
exact; the cost is that the frame operator is no longer exactly $I$. Its first
order change is removed by conditioning the Gaussian on a subspace of
codimension $\le d^2$, and its second-order fluctuation is an average over $n$
rows, of size $t^2d/\sqrt n$. The spectral error of the perturbed frame is
therefore $\eta+O(t^2d/\sqrt n+t^2d^2/n+t^3)$, with $t^2$ a large multiple of $\eta$, which is
at most $\theta t^2$ once that multiple and $B$ are large and $t$ is small; whitening
turns a spectral error $\delta$ into a diagonal error at most $2a\delta$.
(In~\cite{Orig} the remainder was bounded by $O(t^3d)$, which forced
$\eta\le c/d^2$ and a separate block-partition argument; the sharper bound
$O(t^3)$ uses the first-order constraint once more, see \cref{lem:remainder}.)

\paragraph{Seed 2: moderate rows ($d\ge A\log 2n$, \cref{sec:moderate}).}
Perturb in the horizontal tangent space of the Grassmannian (so the frame stays
exactly Parseval) with a Gaussian whose covariance is filtered:
$C_\rho=\rho(I-B)(\rho I+B)^{-1}$, where $B$ is the normalised Gram operator of
the linear diagonal map $\mathcal A$. The filter makes the linear diagonal change
$O(t\sqrt{\rho a\log n/n})$, negligible once $\rho\asymp t^2$ and $d\gg\log n$,
while leaving weak (connecting) directions untouched. The price is a
second-order bias $t^2b_*$, which is not small in $\ell_\infty$. A commutator
identity shows that $b_*=\mathcal Am_*$ for an explicit deterministic tangent
direction $m_*$ with $\fro{m_*}^2\le2a/\rho$ (\cref{lem:commutator,lem:mean}). We
simply add the drift $\frac t2m_*$ to the noise: its first-order effect cancels the
bias exactly, and its other effects are $O(at^2/D)$. All remaining errors are
$\le\theta at^2$ by concentration, which is where $d\ge A\log 2n$ is used. Only
a bounded, deterministic set of rows needs a spectral argument, and second
moments suffice there.

\paragraph{Assembly (\cref{sec:assembly}).}
After complementation and normalisation it suffices to treat equal-norm frames
with $2d\le n$ and spectral error $\eta$. If $d<D$ (a large absolute constant),
the Hamilton--Moitra bound $\eta d(d-1)\le D\eta d$ suffices
(\cref{sec:bounded}). If $d\ge D$ and $n\ge Bd^2$, use Seed~1. If $d\ge D$ and
$n<Bd^2$, then $A\log 2n\le A\log(2Bd^2)\le d$ and Seed~2 applies.

\begin{center}
\small
\begin{tabular}{@{}lll@{}}
\toprule
regime & construction & where\\
\midrule
$d<D$ & radial isotropic position, ordered coordinates (Hamilton--Moitra) & \cref{sec:bounded}\\
$d\ge D,\ n\ge Bd^2$ & row-tangent Gaussian, row normalisation, whitening, balancing & \cref{sec:manyrow}\\
$d\ge D,\ 2d\le n<Bd^2$ & filtered horizontal Gaussian with drift, balancing & \cref{sec:moderate}\\
$d>n/2$ & complement & \cref{sec:assembly}\\
\bottomrule
\end{tabular}
\end{center}

\section{The scaling toolbox}\label{sec:toolbox}

Throughout this section $U\in\R^{n\times d}$ is Parseval, $P=UU^T$,
$p=\diag P$, and $L=L_P=\Diag(p)-P\circ P$.

\subsection{Projections, distances, and the squared-Gram Laplacian}

\begin{lemma}[Frame versus projection distance]\label{lem:align}
Let $U,W\in\R^{n\times d}$ be Parseval with projections $P,Q$, and let
$\sigma_1,\dots,\sigma_d\in[0,1]$ be the singular values of $U^TW$. Then
\[
 \min_{R\in O(d)}\fro{U-WR}^2=2\sum_k(1-\sigma_k)\le
 \fro{P-Q}^2=2\sum_k(1-\sigma_k^2)=2\fro{(I-P)W}^2\le2\fro{U-W}^2 .
\]
Right multiplication by $R\in O(d)$ preserves Parsevalness and every row norm.
\end{lemma}

\begin{proof}
$\fro{U-WR}^2=2d-2\tr(U^TWR)$, maximised over $R$ by the polar factor, giving
$2d-2\sum_k\sigma_k$. Next $\fro{P-Q}^2=2d-2\tr(PQ)=2d-2\fro{U^TW}^2$ and
$\fro{(I-P)W}^2=d-\fro{U^TW}^2$. Finally $1-\sigma\le1-\sigma^2\le2(1-\sigma)$
for $\sigma\in[0,1]$.
\end{proof}

\begin{lemma}[Energy identity]\label{lem:energy}
$L$ is the Laplacian of the weighted graph with weights $P_{ij}^2$ $(i\ne j)$,
$L\succeq0$, $L\one=0$, $L_{I-P}=L_P$, and for $x\in\R^n$, $X=\Diag(x)$,
\[
 x^TLx=\tfrac12\sum_{i,j}P_{ij}^2(x_i-x_j)^2=\fro{(I-P)XU}^2
 =\tfrac12\fro{[X,P]}^2 .
\]
\end{lemma}

\begin{proof}
$\sum_jP_{ij}^2=(P^2)_{ii}=p_i$, so the rows of $\Diag(p)-P\circ P$ sum to zero
and its off-diagonal entries are $-P_{ij}^2$; this gives the first formula. Next
$\fro{(I-P)XU}^2=\tr(XPX(I-P))=\sum_ix_i^2p_i-\sum_{i,j}x_ix_jP_{ij}^2$, and
$[X,P]_{ij}=(x_i-x_j)P_{ij}$. For $I-P$ the off-diagonal weights are again
$P_{ij}^2$.
\end{proof}

\begin{lemma}[Scaled projections]\label{lem:scaled}
For $w\in\R^n_{>0}$ put $D=\Diag(w)$, $M=(U^TD^2U)^{-1/2}$ and $W=DUM$. Then
$W$ is Parseval with column space $D\im U$, $\opn M\le1/\min_iw_i$, and
\[
 \fro{WW^T-P}^2=2\fro{(I-P)W}^2\le\frac{2\,w^TLw}{\min_iw_i^2} .
\]
\end{lemma}

\begin{proof}
$W^TW=M U^TD^2U M=I$ and $U^TD^2U\succeq(\min w_i)^2I$. By \cref{lem:align}
and \cref{lem:energy},
$\fro{(I-P)W}^2\le\opn M^2\fro{(I-P)DU}^2=\opn M^2\,w^TLw$.
\end{proof}

\subsection{The log-determinant potential}

For $s\in\R^n$ let $D_s=\Diag(e^{s})$, let $P(s)$ be the projection onto
$D_s\im U$, and
\[
 F(s)=\tfrac12\log\det(U^TD_s^2U).
\]

\begin{lemma}\label{lem:potential}
$F$ is smooth and convex, $F(s+c\one)=F(s)+cd$, and
\[
 \nabla F(s)=\diag P(s),\qquad \nabla^2F(s)=2L_{P(s)} .
\]
Moreover, if $\norm{s}_\infty\le r$ then
\begin{equation}\label{eq:hessian-comparison}
 e^{-4r}L_P\preceq L_{P(s)}\preceq e^{4r}L_P .
\end{equation}
\end{lemma}

\begin{proof}
Write $G=U^TD_s^2U=\sum_ie^{2s_i}u_iu_i^T$. Then
$\partial_iF=e^{2s_i}u_i^TG^{-1}u_i=P(s)_{ii}$ because
$P(s)=D_sUG^{-1}U^TD_s$, and differentiating once more gives
$\partial_j\partial_iF=2\delta_{ij}P(s)_{ii}-2P(s)_{ij}^2$, i.e.\
$\nabla^2F=2L_{P(s)}\succeq0$. For~\eqref{eq:hessian-comparison} let
$W_s=D_sUM_s$ with $M_s=G^{-1/2}$, so $\opn{M_s}\le e^r$, $\opn{D_s}\le e^r$.
Since $(I-P(s))D_sP=(I-P(s))W_sM_s^{-1}U^T=0$,
\[
 (I-P(s))\Diag(x)W_s=(I-P(s))D_s(I-P)\Diag(x)UM_s ,
\]
and \cref{lem:energy} for $W_s$ gives $x^TL_{P(s)}x\le e^{4r}x^TL_Px$. The
lower bound is the upper bound applied to the Parseval frame $W_s$ and the
scaling $-s$, since $D_{-s}W_s=UM_s$ has column space $\im U$.
\end{proof}

\begin{lemma}[Exact scaling inequalities]\label{lem:scaling-ineq}
For $z\in\R^n_{>0}$ let $r$ be the diagonal of the projection onto
$\Diag(\sqrt z)\im U$. Then, coordinatewise,
\[
 Lz\le z\circ(r-p),\qquad L(z^{-1})\le-z^{-1}\circ(r-p).
\]
\end{lemma}

\begin{proof}
Let $G=U^T\Diag(z)U$. The matrix $(G-z_iI)G^{-1}(G-z_iI)$ is positive
semidefinite; evaluating its quadratic form at $u_i$ and using
$u_i^TGu_i=\sum_jP_{ij}^2z_j$, $\norm{u_i}^2=p_i$, $z_iu_i^TG^{-1}u_i=r_i$ gives
$\sum_jP_{ij}^2z_j-2z_ip_i+z_ir_i\ge0$, which is the first inequality. Let
$V\in\R^{n\times(n-d)}$ be an isometry onto $(\im U)^\perp$. The subspaces
$\Diag(\sqrt z)\im U$ and $\Diag(z^{-1/2})\im V$ are orthogonal complements,
so the projection onto the latter has diagonal $\one-r$, while $VV^T=I-P$ has
diagonal $\one-p$ and Laplacian $L$ (\cref{lem:energy}). The first inequality
for $(V,z^{-1})$ is the second inequality.
\end{proof}

\subsection{Static balancing}

\begin{definition}\label{def:poisson}
A weighted graph Laplacian $L$ on $n$ vertices has \emph{Poisson constant}
$K$ if for every $f\perp\one$ there is $g$ with $Lg=f$ and
$\norm{g}_\infty\le K\norm f_\infty$.
\end{definition}

If $L$ has a finite Poisson constant then the graph is connected (on each
component the image of $L$ sums to zero). The Poisson constant is the
$\ell_\infty\to\ell_\infty$ norm of the pseudo-inverse modulo constants; it is
the quantity that controls a scaling in the box $\norm s_\infty\le1$.

\begin{lemma}[Median barrier]\label{lem:median}
Let $L$ have Poisson constant $K$, let $\beta\ge0$ with $2\beta K<1$, let
$z\in\R^n_{>0}$, let $m$ be a median of $z$ (so that $\{z\le m\}$ and
$\{z\ge m\}$ both have at least $n/2$ elements), and suppose
\[
 (Lz)_i\le\beta z_i\ \text{ when }z_i>m,\qquad
 (Lz^{-1})_i\le\beta z_i^{-1}\ \text{ when }z_i<m .
\]
Then $\max_iz_i/\min_iz_i\le(1-2\beta K)^{-2}$.
\end{lemma}

\begin{proof}
\emph{Claim.} If $|S|\ge n/2$, $y\ge0$, $y\le m$ on $S$ and $(Ly)_i\le\beta y_i$
for $i\notin S$, then $(1-2\beta K)\max y\le m$. Indeed, $f=\one-\frac
n{|S|}\one_S$ has $f\perp\one$ and $\norm f_\infty\le1$; take $Lg=f$ with
$\norm g_\infty\le K$ and $h=g+K\one\in[0,2K]^n$, so $(Lh)_i=1$ for $i\notin
S$. With $M=\max y$, the vector $v=y-m\one-\beta Mh$ is $\le0$ on $S$ and
satisfies $(Lv)_i\le\beta y_i-\beta M\le0$ off $S$. If $\max v>0$, it is attained
at some $i\notin S$, and $(Lv)_i=\sum_jw_{ij}(v_i-v_j)\le0$ forces $v_j=v_i$
for all neighbours $j$; by connectedness the maximum propagates to $S$, a
contradiction. Hence $y\le m+2\beta KM$, and evaluating at the maximum proves
the claim.

Apply the claim to $y=z$, $S=\{z\le m\}$, and to $y=z^{-1}$, $S=\{z\ge m\}$
(with median $m^{-1}$). This gives $(1-2\beta K)\max z\le m$ and
$(1-2\beta K)m\le\min z$.
\end{proof}

\begin{theorem}[Static balancing]\label{thm:balancing}
Suppose $L_P$ has Poisson constant $K$, and $q\in\R^n$ satisfies
$\sum_iq_i=d$, $\norm{q-p}_\infty\le\beta$ and $\beta K\le\frac14$. Then there
is a Parseval $W$ with $\diag(WW^T)=q$ and
\[
 \fro{U-W}^2\le\fro{UU^T-WW^T}^2\le8\norm{q-p}_1 .
\]
$W$ is obtained from $U$ by a row scaling with ratio of scales at most $2$,
whitening, and an orthogonal right factor.
\end{theorem}

\begin{proof}
Let $\Phi(s)=F(s)-\ip qs$, so $\partial_i\Phi=P(s)_{ii}-q_i$, and let $s^*$
minimise $\Phi$ over the cube $[0,1]^n$. Write $r=\diag P(s^*)$. Coordinatewise
optimality gives
\begin{equation}\label{eq:kkt}
 r_i\le q_i\ \text{ if }s^*_i>0,\qquad r_i\ge q_i\ \text{ if }s^*_i<1 .
\end{equation}
Let $z=e^{2s^*}$ with median $m$; then $1\le\min z\le m\le\max z\le e^2$. If
$z_i>m$ then $s_i^*>0$, so by
\cref{lem:scaling-ineq} and~\eqref{eq:kkt}
$(Lz)_i\le z_i(r_i-p_i)\le z_i(q_i-p_i)\le\beta z_i$; if $z_i<m$ then $s_i^*<1$
and likewise $(Lz^{-1})_i\le\beta z_i^{-1}$. By \cref{lem:median},
$\max s^*-\min s^*\le-\log(1-2\beta K)\le\log2<1$. Hence $s^*$ does not meet
both faces: either $s^*_i>0$ for all $i$, and then $r\le q$, or some $s^*_j=0$,
all $s_i^*<1$, and then $r\ge q$. Since $\sum r_i=d=\sum q_i$, in both cases
$r=q$.

Let $w=e^{s^*}$, $D=\Diag(w)$, $m_0=\min w$, $W_0=DU(U^TD^2U)^{-1/2}$, and
$\delta=q-p$. \Cref{lem:scaling-ineq} with $r=q$ gives $Lz\le z\circ\delta$
for $z=w^2$, hence $z^TLz\le\sum_iz_i^2\delta_i\le16m_0^4\norm\delta_1$, using
$z_i\le4m_0^2$. Since $(z_i-z_j)^2\ge4m_0^2(w_i-w_j)^2$,
\cref{lem:energy} gives $w^TLw\le4m_0^2\norm\delta_1$, and
\cref{lem:scaled} gives $\fro{W_0W_0^T-P}^2\le8\norm\delta_1$. Finally take
$W=W_0R$ with $R$ from \cref{lem:align}.
\end{proof}

\begin{remark}
No existence theorem for scalings, no flow, and no genericity is used: the box
supplies a minimiser, the median barrier keeps it away from one face, and the
trace identity $\sum_ir_i=d$ does the rest.
\end{remark}

\subsection{Local correction in the dual norm}

The balancing theorem needs $\ell_\infty$ control of the target. The next lemma
handles a different kind of target error: one that is small in the norm dual
to a Laplacian with a spectral gap, although possibly large in $\ell_\infty$.

\begin{lemma}[Local correction]\label{lem:local}
Let $W$ be Parseval with projection $P_W$ and diagonal $p_W$. Let $J$ be
symmetric with $J\succeq\sigma I$ on $\one^\perp$ ($\sigma>0$) and
\[
 c_1J\preceq L_{P_W}\preceq C_1J\quad\text{on }\one^\perp .
\]
Let $f\perp\one$ satisfy $\ip fx^2\le E\,x^TJx$ for $x\perp\one$, and let
$0<R\le r$ with
\begin{equation}\label{eq:local-small}
 E<\gamma^2\sigma R^2,\qquad \gamma:=2c_1e^{-4r}.
\end{equation}
Then there is $s\perp\one$ with $\norm s_\infty<R$, $s^TJs\le E/\gamma^2$, such
that the Parseval frame $W'$ obtained from $\Diag(e^s)W$ by whitening satisfies
$\diag(W'W'^T)=p_W+f$ and $L_{P_{W'}}\succeq c_1e^{-4r}J$ on $\one^\perp$.
After an orthogonal right factor,
\[
 \fro{W'-W}^2\le\fro{P_{W'}-P_W}^2\le2C_1e^{4r}E/\gamma^2 .
\]
\end{lemma}

\begin{proof}
Let $F$ be the potential of $W$ and $\Phi(s)=F(s)-F(0)-\ip{p_W+f}s$; minimise
$\Phi$ over the compact convex set $\Omega=\{s\perp\one:\norm s_\infty\le R\}$.
For $s\in\Omega$ with $\norm s_\infty=R$, \cref{lem:potential} gives
\[
 \ip{\nabla\Phi(s)}s=\int_0^12s^TL_{P(us)}s\,du-\ip fs
 \ge2c_1e^{-4r}s^TJs-\sqrt{E\,s^TJs}>0,
\]
because $s^TJs\ge\sigma\norm s_2^2\ge\sigma R^2>E/\gamma^2$. So $-s$ is a
feasible descent direction and the minimiser $s$ lies in the relative interior
of $\Omega$. There $\nabla\Phi(s)\in\R\one$, and $\sum_i\nabla\Phi(s)_i=d-d-0=0$,
so $\diag P(s)=p_W+f$. Pairing $\nabla\Phi(s)=0$ with $s$ as above gives
$\gamma\, s^TJs\le\sqrt{E s^TJs}$, i.e.\ $s^TJs\le E/\gamma^2$, and then
$\norm s_\infty\le\norm s_2<R$ by~\eqref{eq:local-small}. The Laplacian bound
is~\eqref{eq:hessian-comparison}. For the distance, \cref{lem:scaled} with
$w=e^s$, $\min w_i\ge e^{-r}$ and $|e^{s_i}-e^{s_j}|\le e^r|s_i-s_j|$ gives
\[
 \fro{P_{W'}-P_W}^2\le2e^{2r}(e^s)^TL_{P_W}e^s\le2e^{4r}s^TL_{P_W}s
 \le2C_1e^{4r}s^TJs ,
\]
and \cref{lem:align} supplies the right factor.
\end{proof}

\subsection{Poisson constants from a dense core}

\begin{lemma}[Dense core with exceptional vertices]\label{lem:core}
Let $L$ be the Laplacian of nonnegative symmetric weights $w_{ij}$ ($w_{ii}=0$)
on $G\sqcup B$, $G\neq\varnothing$, $|B|=b$. Assume
\[
 \sum_{j\in G}\min\{w_{ij},w_{kj}\}\ge\rho>0\quad(i,k\in G),\qquad
 \max_i\sum_jw_{ij}\le D_{\max},
\]
and, if $B\ne\varnothing$, that $C=L_{BB}$ is invertible with
$\norm{C^{-1}}_{\infty\to\infty}\le R$. Then $L$ has Poisson constant
\[
 K=\frac{1+\min\{b,D_{\max}R\}}{\rho}+R
\]
(with $R=b=0$ if $B=\varnothing$).
\end{lemma}

\begin{proof}
Write $L=\begin{pmatrix}A&-E\\-E^T&C\end{pmatrix}$ with $E=(w_{ij})_{i\in
G,j\in B}\ge0$. $C$ is a symmetric $Z$-matrix with $C\one_B=E^T\one_G\ge0$;
since it is invertible and positive semidefinite it is positive definite, and
$C^{-1}\ge0$ entrywise (if $Cx=y\ge0$ then, with $x_-$ the negative part,
$0\le x_-^Ty=x_-^TCx\le-x_-^TCx_-$, so $x_-=0$). Put $H=C^{-1}E^T\ge0$ and
$T=EC^{-1}=H^T$; then $H\one_G=\one_B$ and $\one_G^TT=\one_B^T$. Each row sum of
$T$ is at most $b$ (entries of a column of $T$ are bounded by its column sum
$1$) and at most $D_{\max}R$ (as $T\one_B\le R\,E\one_B$).

The Schur complement $S=A-EC^{-1}E^T$ is the
Laplacian on $G$ with weights $\tilde w_{ik}=w_{ik}+(EC^{-1}E^T)_{ik}\ge w_{ik}$
($i\ne k$)
(its rows sum to zero because $C^{-1}E^T\one_G=\one_B$). Let $f\perp\one$ with
$\norm f_\infty\le1$, and put $h=f_G+Tf_B$; then $\sum h=\sum f=0$ and
$\norm h_\infty\le\Phi:=1+\min\{b,D_{\max}R\}$. The weights $\tilde w$ inherit
the overlap bound, so the graph on $G$ is connected and $Sy=h$ is solvable. If
$i,k\in G$ maximise and minimise $y$,
\[
 h_i-h_k=\sum_{j\in G}\bigl[\tilde w_{ij}(y_i-y_j)+\tilde w_{kj}(y_j-y_k)\bigr]
 \ge\rho\,(y_i-y_k),
\]
so after adding a constant $\norm y_\infty\le\Phi/\rho$. Then $g_G=y$,
$g_B=Hy+C^{-1}f_B$ solves $Lg=f$ (block substitution), and
$\norm{g_B}_\infty\le\Phi/\rho+R$ because $H$ is row-stochastic.
\end{proof}

\begin{corollary}\label{cor:core}
Suppose $b\le n/10$, $\gamma>0$, and every $i\in G$ has at least $\frac45n$
indices $j\ne i$ with $w_{ij}\ge\gamma/n$. Then the overlap hypothesis holds with
$\rho=\gamma/2$. Moreover:
\begin{enumerate}[label=(\alph*)]
\item if $L\succeq\lambda\Pi$ with $\lambda>0$, then $R\le2b/\lambda$;
\item if $w_{ij}=P_{ij}^2$ for a projection $P$, $a>0$, $P_{ii}\ge a/2$ on $B$,
$\sum_{i\in B}P_{ii}\le\frac14$, and $\max_iP_{ii}\le C_2a$, then
$R\le4/a$ and $D_{\max}R\le4C_2$.
\end{enumerate}
\end{corollary}

\begin{proof}
Two sets of at least $\frac45n$ indices share at least $\frac35n$, of which at
least $\frac12n$ lie in $G$, giving overlap $\ge\gamma/2$. (a) For $x$ supported
on $B$, $x^TCx=x^TLx\ge\lambda(\norm x^2-(\sum x_i)^2/n)\ge\lambda(1-b/n)\norm
x^2$, so $|(C^{-1})_{ij}|\le\opn{C^{-1}}\le(\lambda(1-b/n))^{-1}$; sum over a
row. (b) By $P_{ij}^2\le P_{ii}P_{jj}$,
$(C\one_B)_i=P_{ii}-\sum_{j\in B}P_{ij}^2\ge P_{ii}(1-\sum_{j\in B}P_{jj})\ge
a/4$; so $C$ is strictly diagonally dominant, $C^{-1}\ge0$, and $C^{-1}\one_B\le
(4/a)\one_B$. Also $D_{\max}\le\max_iP_{ii}$.
\end{proof}

\subsection{From a seed to an exact frame}

The two constructions produce the following object.

\begin{definition}[Seed]\label{def:seed}
Let $\Theta\ge2$ and $t>0$. A \emph{$(\Theta,t)$-seed} for an $n\times d$ frame
$X$ is a Parseval $V\in\R^{n\times d}$ with
\[
 \fro{V-X}^2\le\Theta t^2d,\qquad
 L_{VV^T}\text{ has Poisson constant }K\le\frac{\Theta}{at^2},\qquad
 \max_i\bigl|(VV^T)_{ii}-a\bigr|\le\frac{at^2}{4\Theta}.
\]
\end{definition}

\begin{corollary}\label{cor:seed}
If $X$ has a $(\Theta,t)$-seed, there is an ENP frame $W$ with
$\fro{X-W}^2\le4\Theta t^2d$.
\end{corollary}

\begin{proof}
Apply \cref{thm:balancing} to $V$ with $q=a\one$:
$\beta K\le\frac14$, and $\fro{V-W}^2\le8n\cdot\frac{at^2}{4\Theta}\le2t^2d$.
Then $\fro{X-W}^2\le2\fro{X-V}^2+2\fro{V-W}^2\le(2\Theta+4)t^2d\le4\Theta
t^2d$.
\end{proof}

\section{Bounded rank: the Hamilton--Moitra estimate}\label{sec:bounded}

This section reproduces the radial-isotropic and ordered-coordinate argument of
Hamilton and Moitra~\cite{HM}, with a direct compactness proof of the existence
of the radial scaling (as in~\cite{Orig}). It is used for $d$ below a large
absolute constant, and it also shows that ENP frames exist.

\begin{lemma}\label{lem:HM}
Let $1\le d\le n$, let $X\in\R^{n\times d}$ have $\norm{x_i}^2=a$ for all $i$,
and let $\eta=\opn{X^TX-I}$. There is an ENP frame $W$ with
$\fro{X-W}^2\le\eta d(d-1)$.
\end{lemma}

\begin{proof}
If $d=1$ then $X^TX=na=1$ and $W=X$ works. Let $d\ge2$.

\emph{Step 1: full spark.} Suppose first that every $d\times d$ minor
$\mu_S=\det(X_S)^2$ ($S\subset[n]$, $|S|=d$) is nonzero. Consider
$F(s)=\frac12\log\det(X^T\Diag(e^{2s})X)-a\sum_is_i$. By Cauchy--Binet,
\[
 F(s)=\tfrac12\log\sum_{|S|=d}\mu_S\,e^{T(s)_S},\qquad
 T(s)_S=2\Bigl(\sum_{i\in S}s_i-a\sum_{i}s_i\Bigr),
\]
where $T:\R^n\to\R^{\binom nd}$ is linear and $\sum_ST(s)_S=0$ because each $i$
lies in a fraction $d/n=a$ of the $d$-sets. The function
$\mathcal E(y)=\sum_S\mu_Se^{y_S}$ restricted to the subspace $\im T$ has the
nonempty sublevel set $\{\mathcal E\le\mathcal E(0)\}$, on which
$y_S\le\log(\mathcal E(0)/\mu_S)$ for every $S$ and hence, by $\sum_Sy_S=0$,
also $y_S$ is bounded below. This set is compact, so $\mathcal E|_{\im T}$
attains its minimum at some $T(s^*)$, and $s^*$ minimises $F$.

Let $D=\Diag(e^{s^*})$, $M=(X^TD^2X)^{-1/2}$ and $W=DXM$. The computation in
the proof of \cref{lem:potential} (which does not use $X^TX=I$) gives
$\partial_iF=\norm{w_i}^2-a$, so $W$ is ENP, and since
$w_i=e^{s^*_i}Mx_i$ has norm $\sqrt a$,
\[
 w_i=\sqrt a\,\frac{Mx_i}{\norm{Mx_i}}.
\]

\emph{Step 2: ordered coordinates.} Replacing $X,W$ by $XO,WO$ for a suitable
$O\in O(d)$ changes neither the hypotheses nor the distance, and makes
$M=\Diag(\lambda_1,\dots,\lambda_d)$ with $0<\lambda_1\le\dots\le\lambda_d$.
Then $w_{ij}=\lambda_jx_{ij}/\sqrt{c_i}$ with
$c_i=a^{-1}\sum_\ell\lambda_\ell^2x_{i\ell}^2$, so $w_{ij}$ has the sign of
$x_{ij}$ and $w_{ij}^2-x_{ij}^2=x_{ij}^2(\lambda_j^2/c_i-1)$ is nonpositive and
then nonnegative as $j$ increases. As $\sum_j(x_{ij}^2-w_{ij}^2)=a-a=0$, the
prefix sums $\Delta_i(k)=\sum_{j\le k}(x_{ij}^2-w_{ij}^2)$ are nonnegative,
vanish at $k=d$, and
\[
 \sum_j|x_{ij}^2-w_{ij}^2|=2\max_k\Delta_i(k)\le2\sum_{k=1}^{d-1}\Delta_i(k).
\]
Since $W^TW=I$ and $(X^TX)_{jj}\le1+\eta$,
$\sum_i\Delta_i(k)=\sum_{j\le k}((X^TX)_{jj}-1)\le k\eta$. Equal signs give
$(x_{ij}-w_{ij})^2\le|x_{ij}^2-w_{ij}^2|$, and therefore
\[
 \fro{X-W}^2\le2\eta\sum_{k=1}^{d-1}k=\eta d(d-1).
\]

\emph{Step 3: general $X$.} Let $G$ be a Vandermonde matrix with rows
$(1,\theta_i,\dots,\theta_i^{d-1})$ for distinct $\theta_i$; all its maximal
minors are nonzero. Each maximal minor of $X+hG$ is a polynomial in $h$ with
nonzero leading coefficient, so for all but finitely many $h$ the matrix $X+hG$
is full spark; choose such $h_k\to0$ (so small that no row of $X+h_kG$
vanishes) and normalise each row of $X+h_kG$ to length $\sqrt a$ (a positive diagonal left factor, which preserves full spark).
The resulting $X_k\to X$ have ENP frames $W_k$ with
$\fro{X_k-W_k}^2\le\eta_kd(d-1)$, $\eta_k\to\eta$. The set of ENP frames is
compact; a limit point $W$ of $(W_k)$ is ENP and satisfies the bound.
\end{proof}

\begin{corollary}\label{cor:exist}
For all $1\le d\le n$ an ENP frame exists. Consequently every equal-norm
$X\in\R^{n\times d}$ is within $\fro{X-W}^2\le2\fro X^2+2\fro W^2=4d$ of an ENP
frame $W$.
\end{corollary}

\begin{proof}
Apply \cref{lem:HM} to $x_i=\sqrt a\,e_1$ for all $i$.
\end{proof}

\section{The many-row seed}\label{sec:manyrow}

\begin{theorem}\label{thm:manyrow}
There are absolute constants $B,c_*,C_*>0$ such that the following holds. Let
$2\le d$, $n\ge Bd^2$, let $X\in\R^{n\times d}$ have $\norm{x_i}^2=a$ for all
$i$, and suppose $\eta:=\opn{X^TX-I}$ satisfies $0<\eta\le c_*$. Then there is
an ENP frame $W$ with $\fro{X-W}^2\le C_*\eta d$.
\end{theorem}

In~\cite{Orig} the same construction is analysed under the stronger hypothesis
$\eta\le c/d^2$; the improvement comes from \cref{lem:remainder}.
Throughout, $S=X^TX$, so $\opn{S-I}\le\eta\le\frac12$, and
$e_i=x_i/\sqrt a$.

\subsection{The noise}

In the Frobenius space $\R^{n\times d}$ consider the subspaces
\[
 \mathcal E=\{Z:\ip{x_i}{z_i}=0\ \forall i\},\qquad
 \mathcal F=\{Z\in\mathcal E:\ X^TZ+Z^TX=0\},
\]
and let $m=\dim\mathcal E-\dim\mathcal F\le\dim\mathrm{Sym}_d\le d^2$. For a
standard Gaussian $g\in\R^{n\times d}$ put
\[
 Z_0=n^{-1/2}\Pi_{\mathcal E}g,\qquad Z=n^{-1/2}\Pi_{\mathcal F}g,\qquad
 Z_\perp=Z_0-Z=n^{-1/2}\Pi_{\mathcal E\ominus\mathcal F}g .
\]
$Z$ and $Z_\perp$ are independent. The rows $z_{0,i}=n^{-1/2}(I-e_ie_i^T)g_i$
of $Z_0$ are independent. Since $\Cov(Z)\preceq\Cov(Z_0)$, each row $z_i$ of
$Z$ is Gaussian with covariance $\Sigma_i\preceq n^{-1}(I-e_ie_i^T)$, and
$\Cov(\mathrm{vec}\,Z)\preceq I/n$. Moreover
\begin{equation}\label{eq:mr-noise}
 \E\fro Z^2=\frac{n(d-1)-m}{n}\le d,\qquad
 \E\fro{Z-Z_0}^2=\frac mn\le\frac{d^2}n .
\end{equation}

\begin{lemma}[Row moments]\label{lem:rowmoments}
Let $r_i=\norm{z_i}^2/a$, $\mu_i=\E r_i$ and $\delta_i=1-\mu_i$. Then
$1/d\le\delta_i\le1$, $\sum_ia\delta_i=1+m/n\le2$, and
\[
 \mathcal A:=a\sum_i\E\bigl[(r_i-1)^2(1+r_i)^2\bigr]\le C_{\mathcal A},
\]
with $C_{\mathcal A}$ absolute.
\end{lemma}

\begin{proof}
$\tilde\Sigma_i=\Sigma_i/a\preceq d^{-1}I$ has trace $\mu_i\le(d-1)/d$, so
$\tr\tilde\Sigma_i^2\le1/d$ and $\tr\tilde\Sigma_i^4\le d^{-3}$. For
$\xi_i=r_i-\mu_i$ the Gaussian moment formulas give $\E\xi_i^2=2\tr\tilde\Sigma_i^2\le
2/d$ and $\E\xi_i^4=12(\tr\tilde\Sigma_i^2)^2+48\tr\tilde\Sigma_i^4\le60/d^2$.
Next $\sum_ia\delta_i=d-\E\fro Z^2=1+m/n$ by~\eqref{eq:mr-noise}. Finally
$(r_i-1)^2\le2\xi_i^2+2\delta_i^2$ and $(1+r_i)^2\le8+2\xi_i^2$, so
\[
 \E(r_i-1)^2(1+r_i)^2\le16\E\xi_i^2+4\E\xi_i^4+16\delta_i^2+4\delta_i^2\E\xi_i^2
 \le\frac{C}d+C\delta_i^2 ,
\]
and $a\sum_i(C/d+C\delta_i^2)\le C+Ca\sum_i\delta_i\le3C$.
\end{proof}

\subsection{Row normalisation and the exact second-order identity}

For $0<t\le1$ set
\[
 v_i=\frac{x_i+tz_i}{\sqrt{1+t^2r_i}},\qquad
 v_{0,i}=\frac{x_i+tz_{0,i}}{\sqrt{1+t^2r_{0,i}}},\qquad
 r_{0,i}=\norm{z_{0,i}}^2/a .
\]
Since $z_i\perp x_i$, $\norm{v_i}^2=\norm{v_{0,i}}^2=a$ exactly. The map
$y\mapsto\sqrt a\,y/\norm y$ is the metric projection onto the closed ball of
radius $\sqrt a$ on $\{\norm y\ge\sqrt a\}$, hence $1$-Lipschitz there. Thus
\begin{equation}\label{eq:mr-contract}
 \fro{V-X}\le t\fro Z,\qquad \fro{V-V_0}\le t\fro{Z-Z_0},\qquad
 \opn V\le\opn{X+tZ},\quad\opn{V_0}\le\opn{X+tZ_0},
\end{equation}
the last two because $V=\Diag(\kappa)(X+tZ)$ with $0<\kappa_i\le1$.

Expanding $1/(1+t^2r_i)=1-t^2c_i$ with
\[
 c_i=\frac{r_i}{1+t^2r_i},
\]
and using $X^TZ+Z^TX=0$, one finds the exact identity
\begin{equation}\label{eq:mr-identity}
 V^TV=S+t^2Q(Z)+R_3+R_4,\qquad Q(Z)=Z^TZ-\sum_ir_ix_ix_i^T,
\end{equation}
\[
 R_3=-t^3\sum_ic_i(x_iz_i^T+z_ix_i^T),\qquad
 R_4=-t^4\sum_ic_i(z_iz_i^T-r_ix_ix_i^T).
\]

\begin{lemma}[Mean and fluctuation of $Q$]\label{lem:Q}
$\opn{\E Q(Z)-(I-S)}\le m/n$ and $\E\fro{Q(Z)-\E Q(Z)}^2\le8d^2/n$.
\end{lemma}

\begin{proof}
For $Z_0$: $\E Z_0^TZ_0=n^{-1}\sum_i(I-e_ie_i^T)=I-S/d$ and $\E r_{0,i}=(d-1)/d$,
so $\E Q(Z_0)=I-S/d-(1-1/d)S=I-S$. As $Z_0=Z+Z_\perp$ with independent centred
summands and $Q$ is quadratic, $\E Q(Z)=\E Q(Z_0)-\E Q(Z_\perp)$, and
$\E Q(Z_\perp)$ is the difference of the positive semidefinite matrices
$\E Z_\perp^TZ_\perp$ and $\sum_i\E(\norm{z_{\perp,i}}^2/a)x_ix_i^T$, both of trace
$\E\fro{Z_\perp}^2=m/n$; so its norm is at most $m/n$.

Entry $(\alpha,\beta)$ of $Q$ is the quadratic form $\sum_iz_i^TB_{\alpha\beta,i}z_i$
with $B_{\alpha\beta,i}=\frac12(E_{\alpha\beta}+E_{\beta\alpha})-(e_i)_\alpha(e_i)_\beta I$.
For a Gaussian vector with covariance $\Sigma\preceq I/n$ and symmetric
coefficient matrix $\mathsf B$, the variance of the quadratic form is
$2\fro{\Sigma^{1/2}\mathsf B\Sigma^{1/2}}^2\le2\fro{\mathsf B}^2/n^2$. Since
$\sum_{\alpha,\beta}\fro{B_{\alpha\beta,i}}^2\le2d^2+2d\le4d^2$ for each $i$,
summing over $\alpha,\beta$ and the $n$ blocks gives the claim.
\end{proof}

\begin{lemma}[Remainder]\label{lem:remainder}
Let $\bar c=1/(1+t^2)$, $\Gamma=\Diag(c_i-\bar c)$ and
$\Lambda=\Diag(c_ir_i-\bar c)$. Then
\[
 \opn{R_3}\le2t^3\opn X\fro{\Gamma Z},\qquad
 \opn{R_4}\le t^4\bigl(\opn Z^2+\opn Z\fro{\Gamma Z}+\opn S+\opn X\fro{\Lambda X}\bigr),
\]
and $\E\fro{\Gamma Z}^2\le\mathcal A$, $\E\fro{\Lambda X}^2\le4\mathcal A$.
\end{lemma}

\begin{proof}
Because $\sum_i(x_iz_i^T+z_ix_i^T)=X^TZ+Z^TX=0$, the constant part $\bar c$ of
the weights drops out of $R_3$:
$R_3=-t^3(X^T\Gamma Z+Z^T\Gamma X)$, which gives the first bound. Next
$\sum_ic_iz_iz_i^T=\bar cZ^TZ+Z^T\Gamma Z$ and $\sum_ic_ir_ix_ix_i^T=\bar cS+X^T\Lambda X$,
with $\opn{Z^T\Gamma Z}\le\opn Z\fro{\Gamma Z}$ and similarly for $\Lambda$, and
$\bar c\le1$. For the moments,
\[
 c_i-\bar c=\frac{r_i-1}{(1+t^2r_i)(1+t^2)},\qquad
 c_ir_i-\bar c=\frac{(r_i-1)(1+r_i+t^2r_i)}{(1+t^2r_i)(1+t^2)},
\]
so $|c_i-\bar c|\le|r_i-1|$ and $|c_ir_i-\bar c|\le2|r_i-1|(1+r_i)$. Hence
$\fro{\Gamma Z}^2=a\sum_i(c_i-\bar c)^2r_i\le a\sum_i(r_i-1)^2(1+r_i)^2$ and
$\fro{\Lambda X}^2=a\sum_i(c_ir_i-\bar c)^2\le4a\sum_i(r_i-1)^2(1+r_i)^2$; take
expectations and use \cref{lem:rowmoments}.
\end{proof}

\begin{remark}
The crude bound $\opn{R_3}\le2t^3\sum_ic_i\norm{x_i}\norm{z_i}\approx2t^3d$ is
what forces $\eta\lesssim d^{-2}$ in~\cite{Orig}. Centring the weights at
$\bar c$, which is allowed by the first-order constraint, replaces the $n$
terms of size $a$ by fluctuations of relative size $d^{-1/2}$, whose
$\ell_2$-sum over the rows is $O(1)$.
\end{remark}

\subsection{A dense reference graph}

\begin{lemma}\label{lem:mr-dense}
There are absolute $h,t_0\in(0,1]$ and $c_1>0$ such that for $t\le t_0$, with
probability at least $1-ne^{-c_1(n-1)}$, every $i$ has at least $0.9(n-1)$
indices $j\ne i$ with
\[
 |\ip{v_{0,i}}{v_{0,j}}|\ge h\,t\sqrt{a/n}.
\]
\end{lemma}

\begin{proof}
Write $z_{0,j}/\sqrt a=d^{-1/2}\gamma_j$ with $\gamma_j\sim N(0,I-e_je_j^T)$
independent. Fix $i$ and condition on row $i$; let $f=v_{0,i}/\sqrt a$, a unit
vector. For $j\ne i$,
\[
 \frac{\ip{v_{0,i}}{v_{0,j}}}{a}=\frac{\nu_j}{N_j},\qquad
 \nu_j=u_j+\tfrac t{\sqrt d}\ip f{\gamma_j},\quad u_j=\ip f{e_j},\quad
 N_j=\bigl(1+\tfrac{t^2}d\norm{\gamma_j}^2\bigr)^{1/2}\ge1,
\]
and the target inequality reads $|\nu_j|/N_j\ge ht/\sqrt d$. By Markov,
$\Prob(N_j>2)\le t^2/3$. On $\{N_j\le2\}$ failure requires
$|\nu_j|<2ht/\sqrt d$. Here $\nu_j\sim N(u_j,t^2(1-u_j^2)/d)$. If $|u_j|\le\frac12$
the variance is at least $3t^2/(4d)$ and the Gaussian density bound gives
probability at most $4h/\sqrt{3\pi/2}\le3h$. If $|u_j|>\frac12$ and $2ht\le\frac14$,
failure requires $\frac t{\sqrt d}|\ip f{\gamma_j}|>\frac14$, of probability at most
$16t^2/d$. Choose $h$ and $t_0$ so that $3h+t_0^2/3+8t_0^2\le\frac1{50}$. The
failure indicators for $j\ne i$ are conditionally independent with
probabilities at most $\frac1{50}$, so by the exponential Markov inequality
(with parameter $1$) more than $0.1(n-1)$ failures have probability at most
$\exp(-(n-1)[0.1-\frac{e-1}{50}])\le e^{-c_1(n-1)}$. Take a union bound over
$i$.
\end{proof}

\subsection{Whitening}

\begin{lemma}\label{lem:whiten}
Let $V$ be equal-norm with $\delta=\opn{V^TV-I}\le\frac12$, $U=V(V^TV)^{-1/2}$ and
$P=UU^T$. Then $|P_{ii}-a|\le2a\delta$,
$\sum_j(P-VV^T)_{ij}^2\le6a\delta^2$ for each $i$, and $\fro{U-V}^2\le d\delta^2$.
\end{lemma}

\begin{proof}
$P_{ii}=v_i^T(V^TV)^{-1}v_i\in[a/(1+\delta),a/(1-\delta)]$. With
$E=(V^TV)^{-1}-I$, $\opn E\le\delta/(1-\delta)\le2\delta$, row $i$ of
$P-VV^T=VEV^T$ has squared norm $v_i^TEV^TVEv_i\le4\delta^2(1+\delta)a$.
Finally the singular values $\sigma_k$ of $V$ satisfy $|\sigma_k-1|\le|\sigma_k^2-1|\le\delta$.
\end{proof}

\subsection{Proof of \texorpdfstring{\cref{thm:manyrow}}{Theorem}}

\emph{Constants.} Fix $h,t_0,c_1$ from \cref{lem:mr-dense}. Put $\kappa=h^2/200$ and
$C_{\rm mix}=125/h^2+3$, then $\Theta=\max\{42,C_{\rm mix}\}$ and
$\theta=\min\{1/(4\Theta),h/60,1/16\}$. Let $C_R=300+100\sqrt{C_{\mathcal A}}$, choose $t_1=\min\{t_0,\theta/(16C_R)\}$, and set
\[
 t^2=4\eta/\theta,\qquad c_*=\theta t_1^2/4,
\]
so that $\eta\le c_*$ gives $t\le t_1$. Finally choose $B$ so large that
$n\ge Bd^2$ implies $d^2/n\le\theta/4$, $\sqrt{160d^2/n}\le\theta/4$,
$\sum_{i\in B_0}2a\le\frac14$ for the set $B_0$ below (which has
$|B_0|\le C'd/\kappa$), $n\ge30$, and $ne^{-c_1(n-1)}\le0.01$.

\emph{The event.} By Markov's inequality, \eqref{eq:mr-noise},
\cref{lem:Q,lem:remainder}, the net bound (\cref{lem:gauss}(c)) for $Z$ and
$Z_0$, and \cref{lem:mr-dense}, with positive probability all of the following
hold:
\[
 \fro Z^2\le20d,\quad \fro{Z-Z_0}^2\le\frac{20d^2}n,\quad
 \fro{Q-\E Q}^2\le\frac{160d^2}n,\quad
 \fro{\Gamma Z}^2\le20\mathcal A,\quad \fro{\Lambda X}^2\le80\mathcal A,
\]
$\opn Z,\opn{Z_0}\le16$, and the conclusion of \cref{lem:mr-dense}. Fix such a
sample.

\emph{Spectral error.} By~\eqref{eq:mr-identity},
$V^TV-I=(1-t^2)(S-I)+t^2(\E Q-(I-S))+t^2(Q-\E Q)+R_3+R_4$. Since
$\opn X\le\sqrt2$, $\opn S\le\frac32$ and $t\le1$, \cref{lem:remainder} and the
event give
$\opn{R_3}+\opn{R_4}\le t^3\bigl(2\sqrt{40\mathcal A}+256+16\sqrt{20\mathcal A}
+\tfrac32+\sqrt{160\mathcal A}\bigr)\le C_Rt^3$, and therefore
\[
 \delta:=\opn{V^TV-I}\le\eta+t^2\frac{d^2}n+t^2\sqrt{\frac{160d^2}n}+C_Rt^3
 \le\frac\theta4t^2+\frac\theta4t^2+\frac\theta4t^2+\frac\theta{16}t^2<\theta t^2 .
\]

\emph{Graph.} By~\eqref{eq:mr-contract}, $\opn V,\opn{V_0}\le\sqrt2+16\le18$
and
$\fro{VV^T-V_0V_0^T}\le36\,t\fro{Z-Z_0}$, so
$\fro{VV^T-V_0V_0^T}^2\le C't^2d^2/n$. Let $B_0$ be the set of rows whose
contribution to this sum exceeds $\kappa at^2$; then $|B_0|\le C'd/\kappa$. A row
$i\notin B_0$ has at most $\kappa at^2/((h/2)^2t^2a/n)=n/50$ indices with
$|(VV^T-V_0V_0^T)_{ij}|\ge\frac h2t\sqrt{a/n}$, so by \cref{lem:mr-dense} it has at
least $0.85n$ indices $j\ne i$ with $|(VV^T)_{ij}|\ge\frac h2t\sqrt{a/n}$.

Let $U=V(V^TV)^{-1/2}$, $P=UU^T$. By \cref{lem:whiten} a row loses at most
$6a\theta^2t^4/((h/4)^2t^2a/n)\le n/20$ of these indices, so each $i\notin
B_0$ has at least $0.8n$ indices $j\ne i$ with $P_{ij}^2\ge(h^2/16)at^2/n$.
Also $a/2\le P_{ii}\le2a$ and $\sum_{i\in B_0}P_{ii}\le\frac14$.
\Cref{lem:core} with $G=B_0^c$, part (i) with $\vartheta=0.3$ and
$\gamma=h^2at^2/16$, and part (iii) (so $R=\frac8{3a}$, $F\le\frac43$) shows that
$H(L_P)\le\frac{7/3}{0.3\gamma}+\frac8{3a}\le C_{\rm mix}/(at^2)$.

\emph{Seed.} $\fro{U-X}^2\le2d\delta^2+2t^2\fro Z^2\le42t^2d$, the diagonal error
is at most $2a\delta\le2\theta at^2\le at^2/(2\Theta)$, and the barrier constant
is at most $\Theta/(at^2)$. So $U$ is a $(\Theta,t)$-seed, and
\cref{cor:seed} gives an ENP frame within $3\Theta t^2d=(12\Theta/\theta)\eta d$
of $X$.
\qed

\section{The moderate-row seed}\label{sec:moderate}

\begin{theorem}\label{thm:moderate}
There are absolute constants $A,C_m,\eps_0>0$ such that the following holds.
Let $U\in\R^{n\times d}$ be Parseval with $2d\le n$, $d\ge A\log(2n)$, and
$|\norm{u_i}^2-a|\le\eps a$ for all $i$, where $0<\eps\le\eps_0$. Then there is
an ENP frame $\widehat U$ with $\fro{\widehat U-U}^2\le C_m\eps d$.
\end{theorem}

Throughout this section $P=UU^T$, $p=\diag P$, $D_p=\Diag(p)$, $L=L_P$, and
$\eps_0\le\frac12$, so $\frac a2\le p_i\le\frac{3a}2\le\frac34$. For a symmetric
$M\in\R^{n\times n}$, $\Lap(M)=\sum_{i<j}M_{ij}^2(e_i-e_j)(e_i-e_j)^T$ is the
Laplacian with weights $M_{ij}^2$; thus $L=\Lap(P)$ and $x^T\Lap(M)x=\frac12\fro{[\Diag
x,M]}^2$. For nonnegative symmetric weights $w$, $\Lap[w]$ denotes the
Laplacian with weights $w_{ij}$. We use $\opn{\Lap[w]}\le2\max_i\sum_jw_{ij}$.

\subsection{The horizontal tangent space and the filtered covariance}

Fix an isometry $V\in\R^{n\times(n-d)}$ onto $(\im U)^\perp$, and let
$\mathcal Z=\R^{(n-d)\times d}$ with the Frobenius inner product. For
$Z\in\mathcal Z$ the matrix $H=VZ$ satisfies $U^TH=0$; its rows are
$h_i=Z^Tv_i$. Define
\[
 A:\mathcal Z\to\R^n,\quad (AZ)_i=v_i^TZu_i=(HU^T)_{ii},\qquad
 A^*x=N_x:=V^T\Diag(x)U .
\]
Then $AA^*=L$, i.e.\ $\fro{N_x}^2=x^TLx$ (\cref{lem:energy}). Let
\[
 \bar A=D_p^{-1/2}A,\qquad S=\bar A\bar A^*=D_p^{-1/2}LD_p^{-1/2},\qquad
 B=\bar A^*\bar A .
\]
Since $L\succeq0$ and $S=I-D_p^{-1/2}(P\circ P)D_p^{-1/2}$ with $P\circ P\succeq0$
(Schur product theorem),
\begin{equation}\label{eq:S}
 0\preceq S\preceq I,\qquad 0\preceq B\preceq I,\qquad
 \tr(I-S)=\sum_i\frac{P_{ii}^2}{p_i}=d .
\end{equation}
Fix an orthonormal eigenbasis of $S$, and for each of its members $y$ with
eigenvalue $\mu>0$ put $f_y=D_p^{-1/2}y$ and
$Z_y=\mu^{-1/2}\bar A^*y=\mu^{-1/2}N_{f_y}$. Since $BZ_y=\mu Z_y$ and
$\ip{Z_y}{Z_{y'}}=\ip y{Sy'}/\sqrt{\mu\mu'}$, the $Z_y$ form an orthonormal
eigenbasis of $B$ on $(\ker B)^\perp=\im\bar A^*$; and since $p_i\ge a/2$,
\begin{equation}\label{eq:finfty}
 \norm{f_y}_\infty\le\sqrt{2/a}.
\end{equation}

\begin{definition}[Filtered covariance]
For $\rho>0$ let $C_\rho=\rho(I-B)(\rho I+B)^{-1}$ on $\mathcal Z$, and let
$Z\sim N(0,C_\rho/n)$, $H=VZ$, $Y=HU^T+UH^T$.
\end{definition}

\begin{lemma}\label{lem:filter}
\begin{enumerate}[label=(\alph*)]
\item $0\preceq C_\rho\preceq I$; $C_\rho=I$ on $\ker B$.
\item $\Cov(AZ)=n^{-1}D_p^{1/2}\,\rho S(I-S)(\rho I+S)^{-1}D_p^{1/2}\preceq
(\rho/n)D_p$.
\item $I-C_\rho=B+R_\rho$ with $R_\rho=\sum_{\mu>0}r_\mu\,Z_y\otimes Z_y$,
$r_\mu=\mu(1-\mu)/(\rho+\mu)\ge0$, and
\[
 \sum_{\mu>0}r_\mu\le d,\qquad \sum_{\mu>0}\frac{r_\mu}{\sqrt\mu}\le\frac d{2\sqrt\rho},
 \qquad\sum_{\mu>0}\frac{r_\mu}\mu\le\frac d\rho .
\]
\end{enumerate}
\end{lemma}

\begin{proof}
(a) The eigenvalues of $C_\rho$ are $\rho(1-\mu)/(\rho+\mu)\in[0,1]$. (b) By
singular value calculus $\bar AC_\rho\bar A^*=\rho S(I-S)(\rho I+S)^{-1}$, whose
eigenvalues $\rho\mu(1-\mu)/(\rho+\mu)$ are at most $\rho$. (c) On the
$\mu$-eigenspace, $1-\rho(1-\mu)/(\rho+\mu)=\mu+\mu(1-\mu)/(\rho+\mu)$. The sums
follow from $\mu/(\rho+\mu)\le1$, $\sqrt\mu/(\rho+\mu)\le1/(2\sqrt\rho)$,
$1/(\rho+\mu)\le1/\rho$, and $\sum_{\mu>0}(1-\mu)\le\tr(I-S)=d$.
\end{proof}

\subsection{The second-order mean}

For $Z\in\mathcal Z$ define the quadratic diagonal
\[
 q_i(Z)=\norm{h_i}^2-\norm{Zu_i}^2=\norm{Z^Tv_i}^2-\norm{Zu_i}^2 .
\]
For a covariance $\Sigma$ on $\mathcal Z$ write $q(\Sigma)=\E q(G)$ with
$G\sim N(0,\Sigma)$; $q(\cdot)$ is linear, and $q(\zeta\otimes\zeta)=q(\zeta)$.

\begin{lemma}[Mean]\label{lem:mean}
$\E q(Z)=a\one-p-b_0-b_*$, where
\[
 b_0=\frac1nL(p^{-1}),\quad \norm{b_0}_\infty\le\frac{12\eps}n,\qquad
 b_*=\frac1n\sum_{\mu>0}r_\mu\,q(Z_y),
\]
and $b_*\perp\one$ with
\begin{equation}\label{eq:bias-dual}
 \ip x{b_*}^2\le\frac{2a}\rho\,x^TLx\qquad(x\in\R^n).
\end{equation}
\end{lemma}

\begin{proof}
For $\Sigma=I$: $\E\norm{Z^Tv_i}^2=d\norm{v_i}^2=d(1-p_i)$ and
$\E\norm{Zu_i}^2=(n-d)p_i$, so $q(I/n)=a\one-p$. Next
$B=\sum_j\zeta_j\otimes\zeta_j$ with $\zeta_j=\bar A^*e_j=v_ju_j^T/\sqrt{p_j}$, and
\[
 q_i(\zeta_j)=\frac{(I-P)_{ij}^2p_j-(1-p_j)P_{ij}^2}{p_j},\qquad
 \sum_jq_i(\zeta_j)=1-p_i-\sum_j\frac{P_{ij}^2}{p_j}+p_i=(Lp^{-1})_i ,
\]
using $\sum_j(I-P)_{ij}^2=1-p_i$ and $\sum_jP_{ij}^2=p_i$. Also
$(Lp^{-1})_i=\sum_jP_{ij}^2(p_i^{-1}-p_j^{-1})$ and
$|p_i^{-1}-p_j^{-1}|\le2\eps a/((1-\eps)a)^2$, so
$|(Lp^{-1})_i|\le(1+\eps)a\cdot8\eps/a\le12\eps$. With \cref{lem:filter}(c)
this gives the formula for $\E q(Z)=q(C_\rho/n)$.

For~\eqref{eq:bias-dual} put $X=\Diag(x)$, $D_x=V^TXV$, $A_x=U^TXU$. Then
$\ip x{q(Z)}=\ip Z{D_xZ-ZA_x}$, and since $X$ and $F=\Diag(f)$ commute,
\begin{equation}\label{eq:commute}
 D_xN_f-N_fA_x=V^T[X(I-P)F-FPX]U=V^T[F(I-P)X-XPF]U=D_fN_x-N_xA_f .
\end{equation}
For $Z=Z_y=\mu^{-1/2}N_f$, $f=f_y$, this gives
$\ip x{q(Z_y)}=\mu^{-1}\ip{N_f}{D_fN_x-N_xA_f}$, and as $\fro{N_f}^2=\mu$,
$\fro{N_x}^2=x^TLx$, $\opn{D_f},\opn{A_f}\le\norm f_\infty$,
\[
 |\ip x{q(Z_y)}|\le\frac{2\norm{f_y}_\infty}{\sqrt\mu}\sqrt{x^TLx}
 \le2\sqrt{\frac2{a\mu}}\sqrt{x^TLx} .
\]
Summing with weights $r_\mu/n$ and using \cref{lem:filter}(c),
$|\ip x{b_*}|\le\frac1n\cdot\frac d{2\sqrt\rho}\cdot2\sqrt{2/a}\sqrt{x^TLx}
=\sqrt{2a/\rho}\sqrt{x^TLx}$. Taking $x=\one$ gives $b_*\perp\one$.
\end{proof}

\begin{remark}
The bias $b_*$ need not be small in $\ell_\infty$: near a direct sum whose
components are joined through a few rows it is concentrated there. What
\eqref{eq:bias-dual} says is that it can be removed by a scaling of
energy $O(at^4/\rho)$ (\cref{lem:local}), which is negligible once
$\rho\asymp t^2$ with a large constant.
\end{remark}

\subsection{The expected graph}

For a direction $\zeta\in\mathcal Z$ write
$Y_\zeta=V\zeta U^T+U\zeta^TV^T$; then $\fro{Y_\zeta}^2=2\fro\zeta^2$, and for a
covariance $\Sigma=\sum_k\lambda_k\psi_k\otimes\psi_k$,
$\E_\Sigma\Lap(Y)=\sum_k\lambda_k\Lap(Y_{\psi_k})$.

\begin{lemma}\label{lem:expected-graph}
On $\one^\perp$,
\[
 \E\Lap(Y)\succeq\frac a4I-\Bigl(\frac2n+\frac8d+\frac8\rho\Bigr)L .
\]
\end{lemma}

\begin{proof}
\emph{Unfiltered noise.} For $Z_0\sim N(0,I/n)$ and $i\ne j$,
$Y_{ij}=v_i^TZ_0u_j+v_j^TZ_0u_i$ has variance
$n^{-1}\fro{v_iu_j^T+v_ju_i^T}^2=n^{-1}(p_i+p_j-2p_ip_j-2P_{ij}^2)$. As
$p\le\frac34$, $p_i+p_j-2p_ip_j\ge\frac14(p_i+p_j)\ge\frac a4$, so
$\E\Lap(Y_{Z_0})\succeq\frac a{4n}\Lap[\one\one^T]-\frac2nL=\frac a4\Pi-\frac2nL$.

\emph{A commutator bound.} For $f\in\R^n$, $F=\Diag(f)$, put
$T_f=Y_{N_f}=(I-P)FP+PF(I-P)=FP+PF-2PFP$. Expanding,
\begin{equation}\label{eq:Tf}
 [X,T_f]=(I-2P)F[X,P]+[X,P]F(I-2P),
\end{equation}
with $I-2P$ orthogonal. Hence $x^T\Lap(T_f)x\le4\norm f_\infty^2\,x^TLx$, and for
$f=e_k$, since $E_k[X,P]$ and $[X,P]E_k$ are the $k$th row and column of $[X,P]$,
$x^T\Lap(T_{e_k})x\le2\sum_jP_{kj}^2(x_k-x_j)^2$.

\emph{Base loss.} $B=\sum_k\zeta_k\otimes\zeta_k$ with $\zeta_k=\bar
A^*e_k=N_{e_k}/\sqrt{p_k}$, so $Y_{\zeta_k}=T_{e_k}/\sqrt{p_k}$ and
$\sum_kx^T\Lap(Y_{\zeta_k})x\le\frac2a\cdot2\sum_{k,j}P_{kj}^2(x_k-x_j)^2=\frac8a\,x^TLx$.
So the base part $B/n$ of the removed covariance removes at most
$\frac8{an}L=\frac8dL$.

\emph{Residual loss.} By~\eqref{eq:finfty}, $\Lap(Y_{Z_y})=\mu^{-1}\Lap(T_{f_y})
\preceq\frac8{a\mu}L$, and by \cref{lem:filter}(c) the residual part $R_\rho/n$
removes at most $\frac8{an}\sum_\mu\frac{r_\mu}\mu L\le\frac8\rho L$.
\end{proof}

\subsection{A good sample}

\begin{lemma}\label{lem:sample}
There are absolute $h\in(0,1)$ and $b_{\max}\in\mathbb N$ and, for each
$\eta'\in(0,\frac1{16}]$, a constant $A_{\eta'}$ such that the following holds
for all $\rho>0$, $0<t\le1$, whenever $d\ge A_{\eta'}\log(2n)$. Put
$J=L+at^2\Pi$. With positive probability:
\begin{enumerate}[label=\textup{(S\arabic*)}]
\item\label{S1} $\opn Z\le16$, $\max_i\norm{h_i}^2\le3a$, $\max_i\norm{Y_i}^2\le800a$;
\item\label{S2} $\max_i|(AZ)_i|\le3\sqrt{\rho a\log(2n)/n}$ and
$\max_i|q_i(Z)-\E q_i(Z)|\le\eta'a$;
\item\label{S3} $\opn{\Lap(Y)-\E\Lap(Y)}\le\eta'a$;
\item\label{S4} $|x^T\mathcal C(Y)x|\le\frac{\eta'}t\,x^TJx$ for $x\perp\one$, where
$\mathcal C(Y)=\sum_{i<j}2P_{ij}Y_{ij}(e_i-e_j)(e_i-e_j)^T$;
\item\label{S5} there is a deterministic set $\mathfrak B$ (depending only on
$U,\rho$) with $|\mathfrak B|\le b_{\max}$ such that every $i\notin\mathfrak B$
has at least $0.95n$ indices $j\ne i$ with $|P_{ij}+tY_{ij}|\ge h\,t\sqrt{a/n}$.
\end{enumerate}
\end{lemma}

\begin{proof}
We show that each item fails with probability at most $\frac1{10}$. Note
$C_\rho\preceq I$, so every linear functional $\ip Z\Psi$ has variance at most
$\fro\Psi^2/n$; we use the tools of \cref{app:gauss}.

\ref{S1}: The net bound (\cref{lem:gauss}(d)) with $\sigma^2=1/n$ gives
$\Prob(\opn Z>16)\le2\cdot5^ne^{-8n}$. The vector $h_i=Z^Tv_i$ has covariance
$\preceq n^{-1}I_d$ and trace $\le a$; \cref{lem:gauss}(b) with $u=2a$ (so
$\min\{u^2/\fro\cdot^2,u/\opn\cdot\}\ge\min\{4d,2d\}$) and a union bound give
$\max\norm{h_i}^2\le3a$ off probability $2ne^{-cd}$. Then
$\norm{Y_i}\le\norm{h_i}+\norm{Zu_i}\le\sqrt{3a}+16\sqrt{3a/2}$.

\ref{S2}: $(AZ)_i$ is Gaussian with variance $\le\rho p_i/n\le3\rho a/(2n)$
(\cref{lem:filter}(b)); a union bound gives the first claim off probability
$2n(2n)^{-3}$. The
form $q_i(Z)=\ip Z{\mathsf M_iZ}$, $\mathsf M_iZ=v_iv_i^TZ-Zu_iu_i^T$, has
$\opn{\mathsf M_i}\le1$ and $\fro{\mathsf M_i}^2\le2(d+np_i^2)\le7d$; after
the covariance factor $(C_\rho/n)^{1/2}$ these become $\le1/n$ and
$\le\sqrt{7d}/n$, and \cref{lem:gauss}(b) with $u=\eta'a$ gives failure
probability $2n\exp(-c\eta'^2d)$.

\ref{S3}: $\Cov(\mathrm{vec}\,Y)\preceq\frac2nI$ (the map $Z\mapsto Y_Z$ is
$\sqrt2$ times an isometry) and $\E\norm{Y_i}^2\le3a$ (for unfiltered noise
$\E\norm{Y_i}^2=a(1-p_i)+(1-a)p_i$, and filtering only decreases it). Apply
\cref{lem:sqgauss}.

\ref{S4}: Let $w_{ij}=J^{-1/2}(e_i-e_j)$, so $\norm{w_{ij}}^2\le2/(at^2)$. The
matrix $t J^{-1/2}\mathcal C(Y)J^{-1/2}=\sum_{i<j}tY_{ij}\cdot2P_{ij}w_{ij}w_{ij}^T$
is a Gaussian series whose coefficients have covariance $\preceq\frac{2t^2}nI$, so
by covariance domination (\cref{lem:gauss}(a)) its matrix variance is at most
\[
 \frac{2t^2}n\sum_{i<j}4P_{ij}^2\norm{w_{ij}}^2w_{ij}w_{ij}^T
 \preceq\frac{16}{an}J^{-1/2}LJ^{-1/2}\preceq\frac{16}dI ,
\]
and \cref{lem:gauss}(a) bounds its norm by $C\sqrt{\log(20n)/d}\le\eta'$ off
probability $\frac1{10}$.

\ref{S5}: \emph{Variance bookkeeping.} By the proof of
\cref{lem:expected-graph}, for $i\ne k$ with $P_{ik}^2\le a/16$ the unfiltered
variance of $Y_{ik}$ is at least $a/(8n)$; at most $24$ indices $k$ have
$P_{ik}^2>a/16$. Row $i$ of $T_{e_j}$ is
$\delta_{ij}e_j^TP+P_{ij}e_j^T-2P_{ij}e_j^TP$, of squared norm at most
$3(\delta_{ij}p_j+P_{ij}^2+4P_{ij}^2p_j)$, so the base part removes from row $i$
total variance $\frac1n\sum_j\norm{e_i^TT_{e_j}}^2/p_j\le\frac{6}{an}\cdot5p_i\le
\frac{45}n$, which exceeds $\frac a{32n}$ at no more than $1440n/d$ positions.
The residual part removes total variance
$\frac1n\sum_\mu r_\mu\fro{Y_{Z_y}}^2\le2a$ from the whole matrix. Fix
$\delta_0=1/6400$ and let $\mathfrak B$ be the set of rows from which it removes
more than $\delta_0a$; then $|\mathfrak B|\le2/\delta_0=:b_{\max}$, and a row
$i\notin\mathfrak B$ loses more than $\frac a{32n}$ to it at no more than
$32\delta_0n$ positions. So for $d\ge10^6$ every $i\notin\mathfrak B$ has at least
$0.99n$ indices $k\ne i$ with $\Var(Y_{ik})\ge\frac a{16n}$.

\emph{Small ball.} Let $\phi:\R\to[0,1]$ be $C^1$, equal to $1$ on $[-h,h]$,
vanishing outside $[-2h,2h]$, with $|\phi'|\le2/h$, and for $i\notin\mathfrak B$ let
$F_i=\frac1n\sum_{k\ne i}\phi\bigl((P_{ik}+tY_{ik})/(t\sqrt{a/n})\bigr)$. By the
density bound for a Gaussian of variance $\ge\frac{at^2}{16n}$,
$\E F_i\le0.01+\frac{4h}{\sqrt{2\pi/16}}\le0.01+7h$. As a function of the row
$(Y_{ik})_k$, $F_i$ has gradient norm at most
$\frac1n\cdot\frac2h\cdot\sqrt{n/a}\cdot\sqrt n=2/(h\sqrt a)$; the row depends on
$Z$ through a linear map of norm $\le2$, and $Z=(C_\rho/n)^{1/2}g$ with $g$
standard. So $F_i$ is $4/(h\sqrt d)$-Lipschitz in $g$, and \cref{lem:gauss}(c) gives
$\Prob(F_i>\E F_i+0.01)\le2\exp(-c\,h^2d)$. Take $h=1/500$, so that
$\E F_i+0.01\le0.04$, and a union bound over $i\notin\mathfrak B$. On the
complement, at most $0.04n$ indices $k\ne i$ have $|P_{ik}+tY_{ik}|<ht\sqrt{a/n}$,
so at least $n-1-0.04n\ge0.95n$ have the reverse inequality.

Finally choose $A_{\eta'}$ so large that $d\ge A_{\eta'}\log(2n)$ implies $d\ge10^6$
and all the failure probabilities above are at most $\frac1{10}$ (there are
five items, so their union has probability $<1$).
\end{proof}

\subsection{Retraction}

Fix a sample as in \cref{lem:sample} and set
\[
 G=I+t^2Z^TZ,\qquad U_t=(U+tH)G^{-1/2},\qquad P_t=U_tU_t^T .
\]
Since $U^TH=0$ and $H^TH=Z^TZ$, $(U+tH)^T(U+tH)=G$ and $U_t$ is Parseval.

\begin{lemma}\label{lem:retraction}
There is an absolute $C_E$ such that for $0<t\le1$:
\begin{enumerate}[label=(\alph*)]
\item $\fro{U_t-U}^2\le C_Et^2d$;
\item every row of $\mathsf E=P_t-P-tY$ satisfies $\sum_j\mathsf E_{ij}^2\le C_Eat^4$;
\item $\diag P_t=p+2tAZ+t^2\E q(Z)+t^2(q(Z)-\E q(Z))+\mathsf r$ with
$\norm{\mathsf r}_\infty\le C_Eat^3$.
\end{enumerate}
\end{lemma}

\begin{proof}
By \ref{S1}, $\opn{G^{-1}-I},\opn{G^{-1/2}-I}\le256t^2$, $\opn{U+tH}\le17$,
$\fro{U+tH}^2\le d(1+256t^2)$ and $\norm{u_i+th_i}\le3\sqrt a$. (a)
$\fro{U_t-U}\le\fro{U+tH}\opn{G^{-1/2}-I}+t\fro Z\le Ct\sqrt d$. (b)
$\mathsf E=(U+tH)(G^{-1}-I)(U+tH)^T+t^2HH^T$, and row $i$ has norm at most
$3\sqrt a\cdot256t^2\cdot17+t^2\sqrt{3a}\cdot16$. (c) Insert
$G^{-1}=I-t^2Z^TZ+t^4(Z^TZ)^2G^{-1}$ into
$(P_t)_{ii}=(u_i+th_i)^TG^{-1}(u_i+th_i)$ and use $u_i^Th_i=(AZ)_i$:
\[
 (P_t)_{ii}=p_i+2t(AZ)_i+t^2q_i(Z)-2t^3u_i^TZ^TZh_i-t^4\norm{Zh_i}^2
 +t^4(u_i+th_i)^T(Z^TZ)^2G^{-1}(u_i+th_i),
\]
and the last three terms are $O(at^3)$ by \ref{S1}.
\end{proof}

\begin{lemma}[Regularised seed]\label{lem:regularised}
There are absolute $t_*\in(0,1]$ and $C_J$ such that if $\eta'\le\frac1{16}$,
$\rho=D^2t^2$ with $D\ge8$, and $0<t\le t_*$, then on $\one^\perp$
\[
 \tfrac1{32}J\preceq L_{P_t}\preceq C_JJ,
\]
and every $i\notin\mathfrak B$ has at least $0.9n$ indices $j\ne i$ with
$|(P_t)_{ij}|\ge\frac h2t\sqrt{a/n}$.
\end{lemma}

\begin{proof}
Entrywise $(P+tY)_{ij}^2=P_{ij}^2+2tP_{ij}Y_{ij}+t^2Y_{ij}^2$, so
$\Lap(P+tY)=L+t\mathcal C(Y)+t^2\Lap(Y)$. On $\one^\perp$, by
\ref{S4}, \ref{S3} and \cref{lem:expected-graph},
\[
 \Lap(P+tY)\succeq L-\eta'J+t^2\Bigl(\frac a4-\eta'a\Bigr)I
 -t^2\Bigl(\frac2n+\frac8d+\frac8{D^2t^2}\Bigr)L
 \succeq\Bigl(\frac78-10t^2-\frac18\Bigr)L+\frac{at^2}8I\succeq\frac18J
\]
for $t^2\le1/40$. Also $\Lap(P+tY)\preceq2L+2t^2\Lap(Y)\preceq2L+3200at^2I$
by \ref{S1}. Now $(P_t)_{ij}^2\ge\frac12(P+tY)_{ij}^2-\mathsf E_{ij}^2$ and
$(P_t)_{ij}^2\le2(P+tY)_{ij}^2+2\mathsf E_{ij}^2$, while
$\Lap(\mathsf E)\preceq2C_Eat^4I$ by \cref{lem:retraction}(b). Hence
$L_{P_t}\succeq\frac1{16}J-2C_Eat^4I\succeq\frac1{32}J$ when $t^2\le1/(64C_E)$,
and $L_{P_t}\preceq(6400+4C_E)J$. Finally, in a row $i\notin\mathfrak B$ at most
$C_Eat^4/(\frac{h^2}4t^2\frac an)=4C_Et^2n/h^2\le0.05n$ indices have
$|\mathsf E_{ij}|\ge\frac h2t\sqrt{a/n}$ once $t^2\le h^2/(80C_E)$; combine
with \ref{S5}.
\end{proof}

\subsection{Removing the bias and balancing}

\begin{proof}[Proof of \cref{thm:moderate}]
\emph{Choice of constants.} Fix $h,b_{\max}$ (\cref{lem:sample}) and
$t_*,C_J,C_E$ (\cref{lem:retraction,lem:regularised}); these are absolute. Put
$C'=1024\,e^3C_J$ and $D=\max\{250,40\sqrt{C'}/h\}$, which fixes $\rho=D^2t^2$.
Put
\[
 C_{\rm core}=\frac{32(1+b_{\max})}{h^2}+64e\,b_{\max},\qquad
 C_\Sigma=2C_E+2C',\qquad \Theta=\max\{2,C_{\rm core},C_\Sigma\},
\]
let $K_0=20\Theta$ and $\eta'=1/(20\Theta)$, and choose $A\ge A_{\eta'}$ so large
that $d\ge A\log 2n$ implies $6D\sqrt{\log(2n)/d}\le1/(20\Theta)$ and
$b_{\max}\le n/10$. Finally choose $\eps_0\le\frac12$ so that $t=\sqrt{K_0\eps}$
satisfies $t\le t_*$, $C_Et\le1/(20\Theta)$ and $12\eps\le1/(20\Theta)$ for
$\eps\le\eps_0$. The order of choices is acyclic, and none depends on
$n,d,\eps$ or $U$. Take a sample as in \cref{lem:sample}.

\emph{Bias correction.} Apply \cref{lem:local} to $W=U_t$ with $J$ as above,
$\sigma=at^2$, $c_1=\frac1{32}$, $C_1=C_J$, $f=t^2b_*$ and $r=R=\frac14$. By
\eqref{eq:bias-dual} and $L\preceq J$, $\ip fx^2\le\frac{2at^4}\rho x^TJx$, so
$E=2at^2/D^2$; with $\gamma=2c_1e^{-1}=1/(16e)$, the condition
$E<\gamma^2\sigma R^2=at^2/(4096e^2)$ holds as $D\ge250$. We obtain a Parseval
$U_*$ with
\[
 \diag P_*=\diag P_t+t^2b_*,\qquad
 \fro{U_*-U_t}^2,\ \fro{P_*-P_t}^2\le C'\frac{at^2}{D^2},\qquad
 L_{P_*}\succeq\frac{at^2}{32e}\Pi\ \text{on }\one^\perp,
\]
using $2C_1e^{4r}E/\gamma^2=2C_Je\cdot\frac{2at^2}{D^2}\cdot256e^2=C'at^2/D^2$. Since $\fro{P_*-P_t}^2\le C'at^2/D^2$, each
row loses at most $C'(at^2/D^2)/(\frac{h^2}{16}t^2\frac an)\le0.1n$ of the
indices of \cref{lem:regularised} when $D\ge40\sqrt{C'}/h$; so every
$i\notin\mathfrak B$ has at least $0.8n$ indices $j\ne i$ with
$(P_*)_{ij}^2\ge\frac{h^2}{16}\frac{at^2}n$.

\emph{Poisson constant.} \Cref{cor:core}(a) with $G=\mathfrak B^c$,
$b\le b_{\max}\le n/10$, $\gamma=h^2at^2/16$ and $\lambda=at^2/(32e)$ gives
$R\le64e\,b_{\max}/(at^2)$ and Poisson constant
$K\le(1+b_{\max})\frac{32}{h^2at^2}+\frac{64eb_{\max}}{at^2}\le\frac{C_{\rm core}}{at^2}$.

\emph{Diagonal error.} By \cref{lem:mean,lem:retraction}(c), the bias cancels
exactly:
\[
 \diag P_*-a\one=(1-t^2)(p-a\one)+2tAZ-t^2b_0+t^2\bigl(q(Z)-\E q(Z)\bigr)+\mathsf r .
\]
By \ref{S2} and $\rho=D^2t^2$,
\[
 \frac{\norm{\diag P_*-a\one}_\infty}{at^2}
 \le\frac1{K_0}+6D\sqrt{\frac{\log2n}d}+\frac{12\eps}d+\eta'+C_Et
 \le\frac5{20\Theta}=\frac1{4\Theta} .
\]

\emph{Seed and cost.} By \cref{lem:retraction}(a) and the bias correction,
$\fro{U_*-U}^2\le2C_Et^2d+2C'at^2/D^2\le C_\Sigma t^2d$ with
$C_\Sigma=2C_E+2C'$. So $U_*$ is a $(\Theta,t)$-seed for $U$, and
\cref{cor:seed} gives an ENP frame $\widehat U$ with
$\fro{\widehat U-U}^2\le4\Theta t^2d=4\Theta K_0\eps d$.
\end{proof}

\section{Assembly}\label{sec:assembly}

We use the two seeds through the following equal-norm formulation.

\begin{proposition}\label{prop:equalrow}
There is an absolute constant $C_0$ such that every $X\in\R^{n\times d}$ with
$1\le d$, $2d\le n$ and $\norm{x_i}^2=a$ for all $i$ admits an ENP frame $W$
with
\[
 \fro{X-W}^2\le C_0\,d\,\opn{X^TX-I}.
\]
\end{proposition}

\begin{proof}
Let $\eta=\opn{X^TX-I}$; if $\eta=0$ take $W=X$. Let $A,\eps_0,C_m$ be the
constants of \cref{thm:moderate} and $B,c_*,C_*$ those of \cref{thm:manyrow},
and fix an integer $D\ge2$ with $A\log(2Bd^2)\le d$ for all $d\ge D$.

\emph{Case $d<D$.} \Cref{lem:HM} gives $\fro{X-W}^2\le\eta d(d-1)\le D\eta d$.

\emph{Case $d\ge D$, $n\ge Bd^2$.} If $\eta\le c_*$ use \cref{thm:manyrow};
otherwise \cref{cor:exist} gives cost $4d\le(4/c_*)\eta d$.

\emph{Case $d\ge D$, $n<Bd^2$.} Then $A\log(2n)\le d$. If $\eta\le\eps_0/4$, let
$U=X(X^TX)^{-1/2}$. The singular values of $X$ lie in
$[\sqrt{1-\eta},\sqrt{1+\eta}]$, so $\fro{X-U}^2\le d\eta^2$, and
$\norm{u_i}^2=x_i^T(X^TX)^{-1}x_i\in[a/(1+\eta),a/(1-\eta)]$, so
$|\norm{u_i}^2-a|\le2\eta a$. \Cref{thm:moderate} with $\eps=2\eta$ gives an ENP
$W$ with $\fro{U-W}^2\le2C_m\eta d$, and
$\fro{X-W}^2\le2d\eta^2+4C_m\eta d$. If $\eta>\eps_0/4$ use the trivial cost $4d$.
\end{proof}

\begin{proof}[Proof of \cref{thm:projection}]
Replacing $(P,Q)$ by $(I-P,I-Q)$ preserves $\beta$, ranks of complements and
distances, so we may assume $d\le n/2$; the case $d=0$ is trivial. Write
$P=UU^T$ with $U^TU=I$, $p_i=\norm{u_i}^2$ and $\delta=\beta/a$. If
$\delta\ge\frac12$, any ENP $W$ (\cref{cor:exist}) gives
$\fro{P-WW^T}^2\le2d\le4n\beta$. Otherwise let $x_i=\sqrt a\,u_i/\norm{u_i}$. Then
$\fro{X-U}^2=\sum_i(\sqrt{p_i}-\sqrt a)^2\le\sum_i(p_i-a)^2/a\le\delta^2d$ and
$X^TX=U^T\Diag(a/p)U$, so $\opn{X^TX-I}\le\max_i|a/p_i-1|\le\delta/(1-\delta)\le
2\delta$. \Cref{prop:equalrow} gives an ENP $W$ with
$\fro{U-W}^2\le2\delta^2d+4C_0\delta d\le(1+4C_0)n\beta$, using $\delta d=n\beta$.
Take $Q=WW^T$ and use \cref{lem:align}.
\end{proof}

\begin{proof}[Proof of \cref{thm:main}]
If $\eps\ge\frac12$, an ENP $W$ gives $\fro{X-W}^2\le2\fro X^2+2d\le6d\le12\eps
d$, since $\fro X^2=\tr X^TX\le(1+\eps)d$. Let $\eps<\frac12$ and
$U=X(X^TX)^{-1/2}$, $P=UU^T$. As above $\fro{X-U}^2\le\eps^2d$ and
$P_{ii}=x_i^T(X^TX)^{-1}x_i\in[\frac{1-\eps}{1+\eps}a,\frac{1+\eps}{1-\eps}a]$, so
$\beta=\max_i|P_{ii}-a|\le4\eps a$. \Cref{thm:projection} gives $Q$ with
$\fro{P-Q}^2\le4C\eps d$, and by \cref{lem:align} there is an isometry $W$
spanning $Q$---hence ENP---with $\fro{U-W}^2\le\fro{P-Q}^2$. Then
$\fro{X-W}^2\le2\eps^2d+8C\eps d$.
\end{proof}

\appendix
\section{Gaussian tools}\label{app:gauss}

\begin{lemma}\label{lem:gauss}
The following hold with absolute constants.
\begin{enumerate}[label=(\alph*)]
\item \emph{(Matrix series.)} Let $X=\sum_kg_kA_k$ with independent standard
Gaussians $g_k$ and symmetric $r\times r$ matrices $A_k$, and
$v=\opn{\sum_kA_k^2}=\opn{\E X^2}$. Then $\E\opn X\le C\sqrt{v\log(2r)}$ and
$\Prob\{\opn X>C\sqrt{v\log(2r/\delta)}\}\le\delta$ for $0<\delta\le1$. If instead
$X=\sum_\alpha\xi_\alpha M_\alpha$ with $(\xi_\alpha)$ a centred Gaussian vector of
covariance $\preceq\gamma I$, then $X$ has this form with
$\sum_kA_k^2=\E X^2\preceq\gamma\sum_\alpha M_\alpha^2$.
\item \emph{(Quadratic forms.)} For $g$ standard Gaussian and $\mathsf A$
symmetric, $\Var(g^T\mathsf Ag)=2\fro{\mathsf A}^2$ and
\[
 \Prob\{|g^T\mathsf Ag-\tr\mathsf A|>u\}\le2\exp\Bigl[-\tfrac18\min\Bigl\{\frac{u^2}{\fro{\mathsf A}^2},\frac u{\opn{\mathsf A}}\Bigr\}\Bigr].
\]
For a Gaussian vector $\zeta=\Sigma^{1/2}g$, apply this to
$\Sigma^{1/2}\mathsf A\Sigma^{1/2}$, whose norms are at most $\opn\Sigma$ times those of
$\mathsf A$.
\item \emph{(Lipschitz functions.)} If $F$ is $C^1$ with $\norm{\nabla F}\le\ell$
and $g$ is standard Gaussian, then
$\Prob\{|F(g)-\E F(g)|>u\}\le2\exp(-2u^2/(\pi^2\ell^2))$.
\item \emph{(Nets.)} If $Z$ is an $r\times s$ centred Gaussian matrix with
$\Cov(\mathrm{vec}\,Z)\preceq\sigma^2I$, then
$\Prob\{\opn Z>u\}\le2\cdot5^{r+s}\exp(-u^2/(32\sigma^2))$.
\end{enumerate}
\end{lemma}

\begin{proof}
(a) Let $M_p=\E\tr X^{2p}$. Gaussian integration by parts gives
$M_p=\sum_k\sum_{j=0}^{2p-2}\E\tr(A_kX^jA_kX^{2p-2-j})$. For symmetric $A$, even
$q$ and $0\le j\le q$, diagonalising $X=\sum_l\lambda_le_le_l^T$ and using weighted
AM--GM, $|\tr(AX^jAX^{q-j})|=|\sum_{l,m}A_{lm}^2\lambda_m^j\lambda_l^{q-j}|\le
\sum_{l,m}A_{lm}^2|\lambda_l|^q=\tr(A^2|X|^q)$. Hence $M_p\le(2p-1)vM_{p-1}$ and
$M_p\le r(2p-1)!!\,v^p\le r(2pv)^p$. Markov's inequality for $\tr X^{2p}\ge\opn X^{2p}$
with $p=\lceil\log(r/\delta)\rceil$ and $u^2=2e^2pv$ gives the tail; $p=\lceil\log
2r\rceil$ and Jensen give the mean. For the second statement write
$\xi=\Gamma^{1/2}g$; then $X=\sum_kg_kA_k$ with $A_k=\sum_\alpha(\Gamma^{1/2})_{\alpha
k}M_\alpha$, and for a vector $w$, $w^T\E X^2w=\E\norm{Xw}^2=\tr(\Gamma\mathsf G_w)\le
\gamma\tr\mathsf G_w=\gamma\sum_\alpha\norm{M_\alpha w}^2$, where $\mathsf G_w$ is the
Gram matrix of $(M_\alpha w)_\alpha$.

(b) Diagonalise $\mathsf A$ with eigenvalues $\lambda_j$. For $|s|\le1/(4\opn{\mathsf A})$,
$\log\E e^{s(g^T\mathsf Ag-\tr\mathsf A)}=\sum_j[-s\lambda_j-\frac12\log(1-2s\lambda_j)]
\le2s^2\fro{\mathsf A}^2$, since $-x-\frac12\log(1-2x)=\sum_{k\ge2}(2x)^k/(2k)\le2x^2$
for $|x|\le\frac14$. The exponential Markov inequality with
$s=\min\{u/(4\fro{\mathsf A}^2),1/(4\opn{\mathsf A})\}$, applied to both signs,
gives the tail. The variance follows from $\E g_j^4=3$.

(c) Let $g'$ be an independent copy, $g_\theta=g\cos\theta+g'\sin\theta$ and
$\dot g_\theta=-g\sin\theta+g'\cos\theta$; for each $\theta$ the pair
$(g_\theta,\dot g_\theta)$ is standard Gaussian. By the fundamental theorem of
calculus and Jensen on $[0,\pi/2]$,
\[
 \E e^{s(F(g')-F(g))}\le\frac2\pi\int_0^{\pi/2}\E\exp\Bigl[\frac{\pi s}2\ip{\nabla
 F(g_\theta)}{\dot g_\theta}\Bigr]d\theta\le e^{\pi^2s^2\ell^2/8}.
\]
Jensen in $g'$ gives $\E e^{s(F(g)-\E F)}\le e^{\pi^2s^2\ell^2/8}$; optimise in $s$.

(d) Let $\mathcal N_1,\mathcal N_2$ be $\frac12$-nets of the unit spheres of $\R^r$
and $\R^s$ with at most $5^r$, $5^s$ points. For a vector $w$, choosing a net point
$y'$ near the unit vector $w/\norm w$ gives $\ip w{y'}\ge\frac12\norm w$; applying this
twice gives $\opn Z\le4\max_{x\in\mathcal N_1,y\in\mathcal N_2}|x^TZy|$. Each
$x^TZy$ is Gaussian with variance $\le\sigma^2$; a union bound finishes.
\end{proof}

\begin{lemma}[Squared Gaussian graphs]\label{lem:sqgauss}
Let $Y$ be a centred symmetric Gaussian $n\times n$ matrix whose entries
$(Y_{ij})_{i\le j}$ have covariance $\preceq C_0I/n$, and with
$\max_i\E\sum_jY_{ij}^2\le C_0a$, where $d=an\ge1$. For every $\zeta,\delta\in(0,1)$
there is $C_{\zeta,\delta}$ (depending also on $C_0$) such that if
$d\ge C_{\zeta,\delta}\log(2n)$ then $\opn{\Lap(Y)-\E\Lap(Y)}\le\zeta a$ with
probability at least $1-\delta$.
\end{lemma}

\begin{proof}
$\Lap(Y)=\Diag(\deg)-O(Y)$, where $\deg_i=\sum_{j\ne i}Y_{ij}^2$ and $O(Y)$ is the
off-diagonal part of $Y\circ Y$. Each $\deg_i$ is a Gaussian quadratic form with
coefficient matrix $\mathsf A\succeq0$, $\opn{\mathsf A}\le C_0/n$,
$\tr\mathsf A\le C_0a$, $\fro{\mathsf A}^2\le C_0^2a/n$; by
\cref{lem:gauss}(b) and a union bound, $\max_i|\deg_i-\E\deg_i|\le\zeta a/2$ off
probability $2ne^{-c_\zeta d}$, and also $\E\max_i\deg_i\le C(a+\log(2n)/n)$.

For the off-diagonal part let $Y'$ be an independent copy. By Jensen,
$\E\opn{O(Y)-\E O(Y)}\le\E\opn{O(Y)-O(Y')}$, and $O(Y)-O(Y')$ is the
off-diagonal part of $(Y-Y')\circ(Y+Y')=2\tilde Y\circ\tilde Y'$, where
$\tilde Y=(Y-Y')/\sqrt2$ and $\tilde Y'=(Y+Y')/\sqrt2$ are independent copies of
$Y$. Conditionally on $\tilde Y$, the off-diagonal part of $\tilde Y\circ\tilde Y'$
is the Gaussian series $\sum_{i<j}\tilde Y'_{ij}\tilde Y_{ij}(E_{ij}+E_{ji})$,
whose variance is, by \cref{lem:gauss}(a),
$\preceq\frac{C_0}n\sum_{i<j}\tilde Y_{ij}^2(E_{ii}+E_{jj})\preceq\frac{C_0}n
\max_i\deg_i(\tilde Y)\,I$. Hence
\[
 \E\opn{O(Y)-\E O(Y)}\le C\sqrt{\frac{\log2n}n}\,\E\sqrt{\max_i\deg_i(\tilde Y)}
 \le Ca\Bigl(\sqrt{\frac{\log2n}d}+\frac{\log2n}d\Bigr),
\]
and Markov's inequality finishes the proof.
\end{proof}

\section{Relation to the original manuscript and its formalisation}\label{app:notes}

The construction and almost all estimates are those of~\cite{Orig}. The
differences are as follows.

\begin{enumerate}[label=\arabic*.]
\item \emph{Many-row seed for all small errors.} In~\cite{Orig} the many-row
theorem assumes $\eta\le c/d^2$, because the cubic remainder of the
normalisation identity is bounded termwise by $O(t^3d)$, and the operator norm of
the noise by $t^2\fro{Z}^2\le1$. \Cref{lem:remainder} uses the first-order
constraint $X^TZ+Z^TX=0$ a second time to centre the weights, giving $O(t^3)$,
and the net bound controls $\opn Z$. \Cref{thm:manyrow} therefore holds for
$\eta\le c_*$.
\item \emph{No block partition.} With 1., the global argument needs no random
partition into blocks of size $\lceil Ld^6\rceil$ (and no uniform-subset variance
identity or averaging over block errors): the three regimes $d<D$, $n\ge Bd^2$,
$n<Bd^2$ of \cref{prop:equalrow} cover everything, since $A\log(2Bd^2)\le d$ for
$d\ge D$.
\item \emph{One toolbox.} The distance estimates, the potential, the static
balancing theorem, the local correction lemma and the dense-core Poisson
estimate are stated once (\cref{sec:toolbox}), and both seeds are routed
through \cref{def:seed} and \cref{cor:seed}.
\item \emph{Complete sampling estimates.} The probabilistic statements used by
the moderate seed are collected in \cref{lem:sample} and proved in full,
including the variance bookkeeping behind the dense core.
\end{enumerate}

\paragraph{Formal status.} The Lean~4 development accompanying~\cite{Orig}
proves \cref{thm:main,thm:projection} by the original route (with the block
partition and the $\eta\le c/d^2$ many-row theorem). The argument of
\cref{sec:manyrow,sec:assembly} here is not formalised; the other sections follow
the formalised proof up to constants and presentation.


\begin{thebibliography}{9}
\bibitem{CC}
J.~Cahill and P.~G.~Casazza.
\emph{The Paulsen problem in operator theory}.
Operators and Matrices 7 (2013).
\href{https://arxiv.org/abs/1102.2344}{arXiv:1102.2344}.
\bibitem{HM}
L.~Hamilton and A.~Moitra.
\emph{The Paulsen problem made simple}.
ITCS 2019; Israel J.\ Math.\ (2021).
\href{https://arxiv.org/abs/1809.04726}{arXiv:1809.04726}.
\bibitem{KLLR}
T.~C.~Kwok, L.~C.~Lau, Y.~T.~Lee and A.~Ramachandran.
\emph{The Paulsen problem, continuous operator scaling, and smoothed analysis}.
STOC 2018.
\href{https://arxiv.org/abs/1710.02587}{arXiv:1710.02587}.
\bibitem{LR}
L.~C.~Lau and A.~Ramachandran.
\emph{Optimal bounds for Tyler's M-estimator for elliptical distributions},
Section~7.
\href{https://arxiv.org/abs/2510.13751}{arXiv:2510.13751}, 2025.
\bibitem{Orig}
\emph{A static Gaussian proof of the linear Paulsen bound}.
Manuscript and Lean~4 formalisation (September 2026), produced with an AI
system from a strategy of K.~H\"anni and H.~Eberhard.
\end{thebibliography}
\end{document}
